% !TEX program = xelatex
\documentclass[UTF8,zihao=-4,openany]{ctexbook}

\usepackage[b5paper,top=2.2cm,bottom=2.0cm,left=1.9cm,right=1.7cm,headheight=15pt]{geometry}
\usepackage{amsmath,amssymb}
\usepackage{graphicx}
\usepackage{xcolor}
\usepackage{booktabs,tabularx,array,float,longtable}
\usepackage{enumitem}
\usepackage{hyperref}

\definecolor{titleblue}{HTML}{1F4E79}
\definecolor{deepgreen}{HTML}{2E6F5E}
\definecolor{softgray}{HTML}{F6F7F9}
\definecolor{linegray}{HTML}{D8DEE8}
\definecolor{notebg}{HTML}{FFF8E6}
\definecolor{noteframe}{HTML}{D79A2B}

\hypersetup{
  colorlinks=true,
  linkcolor=titleblue,
  urlcolor=deepgreen,
  citecolor=titleblue,
  pdfauthor={Storm},
  pdftitle={西安交通大学少年班 江苏省苏州中学预科基地 预科一化学笔记}
}

\setlength{\parindent}{2em}
\setlength{\parskip}{0.35em}
\linespread{1.18}
\setlist[itemize]{leftmargin=2.2em,itemsep=0.25em,topsep=0.25em}
\newcolumntype{Y}{>{\centering\arraybackslash}X}
\providecommand{\xrightleftharpoons}[2][]{\mathrel{\mathop{\rightleftharpoons}\limits^{#2}_{#1}}}
\newcommand{\coverdate}{\number\year 年 \number\month 月}

\ctexset{
  chapter = {
    format = \Huge\bfseries\color{titleblue},
    name = {第,章},
    number = \chinese{chapter},
    beforeskip = 0pt,
    afterskip = 2.6em,
  },
  section = {
    format = \Large\bfseries\color{deepgreen},
  },
  subsection = {
    format = \large\bfseries\color{titleblue},
  },
  subsubsection = {
    format = \normalsize\bfseries\color{deepgreen},
  },
}
\pagestyle{headings}

\newenvironment{notebox}
  {\par\smallskip\noindent\begingroup\begin{minipage}{\linewidth}\color{noteframe}\hrule height 0.6pt\smallskip\color{black}\itshape}
  {\par\smallskip\color{noteframe}\hrule height 0.6pt\end{minipage}\endgroup\par\smallskip}

\renewcommand{\contentsname}{目录}
\renewcommand{\chaptername}{第\thechapter 章}

\begin{document}

\frontmatter
\begin{titlepage}
\centering
\vspace*{1.55cm}
{\zihao{2}\songti\bfseries 西安交通大学少年班\par}
\vspace{0.55cm}
{\zihao{2}\songti\bfseries 江苏省苏州中学预科基地\par}
\vspace{1.05cm}
\rule{0.72\linewidth}{1.2pt}\par
\vspace{0.95cm}
{\fontsize{42}{52}\selectfont\songti\bfseries 预科一化学笔记\par}
\vspace{0.65cm}
{\zihao{3}\rmfamily\scshape Chemistry Notes\par}
\vspace{0.95cm}
\rule{0.72\linewidth}{1.2pt}\par
\vfill
{\zihao{4} by Storm\par}
\vspace{0.75cm}
{\zihao{4} \coverdate\par}
\vspace*{1.5cm}
\end{titlepage}

\pagestyle{plain}
\tableofcontents
\cleardoublepage
\pagestyle{headings}
\mainmatter



\chapter{基础知识}
\section{物质的分类}
\subsection{分类方法}
\begin{itemize}
\item 按状态：固（$s$）、液（$l$）、气（$g$）
\end{itemize}
\begin{notebox}
注：不是所有物质都有三态（eg. $CaCO_3$ 只有固态），不是所有物质都只有三态\par
\end{notebox}
\begin{itemize}
\item 按导电性：导体、半导体（eg. $Si$），绝缘体
\end{itemize}

\begin{itemize}
\item 按溶解能力（仅限固体）：
\end{itemize}

\begin{table}[H]
\centering
\renewcommand{\arraystretch}{1.22}
\begin{tabularx}{\linewidth}{Y Y Y Y}
\toprule
难溶 & 微溶 & 可溶 & 易溶 \\
\midrule
$<0.01$g & $0.01$g\textasciitilde{}$1$g & $1$g\textasciitilde{}$10$g & $>10$g \\
\bottomrule
\end{tabularx}
\end{table}

\begin{itemize}
\item 按组成：混合物、纯净物等
\end{itemize}
\subsection{单质与同素异形体}
（1）$C$：$C_{60}$、$C_{100}$、碳纳米管\par

（2）$S$：$S_2$（气体）、$S_6$（固体）、$S_8$（固体）（即 $S$ 的三态变化是化学变化）\par
\subsection{酸与碱}
（1）定义：在水溶液中解离出的阳离子只有 $H^+$ 的物质是酸，在水溶液中解离出的阴离子只有 $OH^-$ 的物质是碱\par

（2）分类：\par
\begin{itemize}
\item 一元酸 / 碱、二元酸 / 碱等
\end{itemize}

\begin{itemize}
\item 可溶酸 / 碱、不溶酸 / 碱（eg. $H_2SiO_3$ 是不溶酸）
\item 强酸（$HClO_4$、$H_2SO_4$、$HCl$、$HBr$、$HI$、$HNO_3$）、弱酸
\item 强碱（$KOH$、$NaOH$、$Ba(OH)_2$、$Ca(OH)_2$）、弱碱
\end{itemize}
\begin{notebox}
注：可溶性碱是强碱加上 $NH_4OH$\par
\end{notebox}
\subsection{盐}
（1）分类：\par
\begin{itemize}
\item 正盐：酸碱中和时正好完全反应产生的盐
\end{itemize}

\begin{itemize}
\item 酸式盐：酸碱中和时酸过量产生的盐（eg. $NaHCO_3$）
\item 碱式盐：酸碱中和时碱过量产生的盐（eg. $Mg(OH)Cl$）
\end{itemize}

（2）改写：\par

$Mg(OH)Cl \to Mg(OH)_2\cdot MgCl_2$\par

$NaHCO_3 \to Na_2CO_3\cdot H_2CO_3$\par

$Cu_2(OH)_2CO_3 \to Cu(OH)_2\cdot CuCO_3$\par
\subsection{两性氢氧化物 $Al(OH)_3$}
$Al(OH)_3\rightleftharpoons Al^{3+}+3OH^-$，体现其“碱”性\par

$Al(OH)_3+OH^-\rightleftharpoons [Al(OH)_4]^-$，体现其“酸”性\par
\begin{notebox}
扩展：$Be(OH)_2$、$Zn(OH)_2$也是两性氢氧化物，且 $Be$ 和 $Al$ 体现了对角线法则\par
\end{notebox}
\subsection{氧化物}
（1）定义：二元含氧，氧为负价（因此 $OF_2$ 为氟化物，而非氧化物）\par

（2）分类：酸性、碱性、两性、不成盐、其它\par

\begin{itemize}
\item 酸性氧化物：与 $OH^-$ 反应生成盐和水（eg. $P_2O_5+6NaOH =\!=\!=2Na_3PO_4+3H_2O $），含氧酸与酸性氧化物存在对应关系（eg. $SO_2$ 对应 $H_2SO_3$）
\end{itemize}

\begin{itemize}
\item 碱性氧化物：与 $H^+$ 反应生成盐和水（eg. $FeO+2HCl =\!=\!= FeCl_2+2H_2O$），碱与碱性氧化物存在对应关系
\end{itemize}

\begin{itemize}
\item 两性氧化物：$Al_2O_3$、$BeO$、$ZnO$（eg. $Al_2O_3+6H^+ =\!=\!= 2Al^{3+}+3H_2O$，$Al_2O_3+2OH^-+3H_2O=\!=\!= 2[Al(OH)_4]^-$）
\end{itemize}

\begin{itemize}
\item 不成盐氧化物（eg. $CO$、$NO$）
\end{itemize}

\begin{itemize}
\item 其它：水和特殊氧化物（eg. $H_2O_2$、$Na_2O_2$、$KO_2$、$KO_3$（臭氧化钾））
\end{itemize}

\section{氧化还原反应（O-R）}
\subsection{定义}
有 $e^-$ 转移的反应 $\to$ 本质\par

\subsection{判断方法}
化合价改变 $\to$ 表象\par

\subsection{四种基本反应类型和 O-R 的关系}
化合、置换、分解、复分解与 O-R 都有交集\par

\begin{notebox}
注：置换与 O-R 的交集：$4Te+5Ni(CO)_4 =\!=\!= 5Ni+ 4Te(CO)_5$，其中 $Te$ 和 $Ni$ 都是 $0$ 价；复分解与 O-R 的交集：$Na_2S+CuSO_4 =\!=\!= CuS\downarrow +Na_2SO_4$，其中 $CuS$ 的价态较复杂，有化合价变化\par
\end{notebox}
\subsection{氧化与还原}
化合价升高：失 $e^-$，被氧化，是还原剂、氧化产物，“升失氧”\par

化合价降低：得 $e^-$，被还原，是氧化剂、还原产物，“降得还”\par

（eg. $Cl_2+H_2O \rightleftharpoons HCl+HClO$，这是一个歧化反应，$Cl_2$ 既是氧化剂又是还原剂，$HCl$ 是还原产物，$HClO$ 是氧化产物）\par
\subsection{氧化性、还原性}
同一物质高价有氧化性、低价有还原性；金属单质多有还原性，非金属单质大多二者都有（在单质中，$F_2$ 只有氧化性，且其氧化性是单质中最强的）\par
\begin{notebox}
扩展：氧化性第二强的单质是 $O_3$\par
\end{notebox}
\subsection{氧化、还原能力}
\begin{itemize}
\item 氧化剂的氧化能力 $>$ 氧化产物，还原剂的还原能力 $>$ 还原产物（特例：电解和可逆反应）
\end{itemize}

\begin{itemize}
\item 金属活动性由弱到强，还原性由弱到强、氧化性由强到弱
\end{itemize}

\begin{itemize}
\item 元素周期表中从上到下单质还原性增强，从左到右单质还原性减弱
\end{itemize}

\subsection{单、双线桥}
单线桥中，箭头方向表示 $e^-$ 移动的方向；双线桥中，箭头由反应物指向生成物，且得失电子数相同。\par
\subsection{化合价升降法配平}
1. 找出氧化剂、还原剂\par

2. 点清升降化合价\par

3. 使得升降化合价相等\par

4. 补齐成分\par

一些例子：\par

$2KMnO_4+16HCl=\!=\!= 2KCl+2MnCl_2+5Cl_2 \uparrow +8H_2O$\par

$4NH_3+5O_2 =\!=\!= 4NO+6H_2O$\par

$Fe_3C+22HNO_3 =\!=\!= 3Fe(NO_3)_3+CO_2\uparrow+13NO_2\uparrow+11H_2O$\par

$FeC_3+18HNO_3=Fe(NO_3)_3+3CO_2\uparrow+15NO_2\uparrow+9H_2O$\par
\begin{notebox}
扩展：$Fe_3C$ 和 $FeC_3$ 中各元素都是 $0$ 价，因为它们是合金。因此有时混合物也可以写化合物，只要它的成分确定（eg. $98\%$ 浓硫酸：$H_2SO_4\cdot 9H_2O$）\par
\end{notebox}
\section{离子方程式}
\subsection{离子反应定义}
有离子参加或生成的反应\par
\subsection{可以拆的物质}
\begin{itemize}
\item 强酸、强碱（$Ca(OH)_2$ 在溶液中拆，在浊液中不拆）
\end{itemize}

\begin{itemize}
\item 可溶性盐（溶不溶视具体情况而定，eg. 侯氏制碱法中 $Na^++NH_3+H_2O+CO_2=\!=\!=NaHCO_3\downarrow+NH_4^+$ 的 $NaHCO_3$ 不溶）
\end{itemize}
\subsection{常见产物}
$SO_2\to SO_4^{2-}$，$S^{2-}\to S$，$CO\to CO_2$，$H_2O_2\to O_2$，$Fe^{2+}\to Fe^{3+}$，$C\to CO_2$，$NH_3\to N_2O/N_2/NO$\par
\begin{notebox}
注：$NH_3$ 反应时催化剂不同，所得产物也不同，这说明了催化剂会影响产物（eg. $4NH_3+5O_2\stackrel{Pt}{=\!=\!=}4NO+6H_2O$，\\[0.25em]
\hspace*{2em}$4NH_3+3O_2=\!=\!=2N_2+6H_2O$）\par
\end{notebox}
\section{物质的量}
\subsection{单位}
摩尔（mol），即 $12.00$g $C-12$ 所含原子的个数；\par

阿伏伽德罗常数的值 $N_A\approx 6.02\times10^{23} $\par
\subsection{转化}
$n\times N_A=N$，$n\times M=m$，$\dfrac{n}{V_c}=c$，$n\times V_m=V$\par

其中，$M$ 在数值上等于相对原子（分子）质量；在标况（STP，即$273.15$K、$1$atm）下，$V_m\approx 22.4$L，即 $1$mol 任何气体在 STP 下的体积都为 $22.4$L\par
\begin{notebox}
注：$22.4$ L 的体积将近是一个棱长 $28$ cm 的正方体或一个半径 $17.5$ cm 的球（这个球略大于篮球）\par
\end{notebox}
\subsection{理想气体状态方程}
$pV=nRT$，其中 $R$ 为常数，$T$ 为开氏温度\par

\subsection{物质的量浓度}
符号为 $c$，$c=\dfrac{n}{V}=\dfrac{1000\rho\eta}{M}$，其中 $\eta$ 是浓度（质量分数），$\rho$ 单位是 g/mL\par
\subsection{配溶液}
\begin{itemize}
\item 步骤：计算 $\to$ 称量/量取 $\to$ 溶解/稀释 $\to$ 转移 $\to$ 洗涤 $\to$ 定容 $\to$ 摇匀 $\to$ 贴标签
\end{itemize}
\begin{notebox}
注：“转移”时用到的玻璃仪器有玻璃棒、烧杯、（）mL 容量瓶，且容量瓶的容积要上取整（eg. 制取 $480$mL 溶液需要的是 $500$mL 容量瓶）\par
\end{notebox}
\begin{itemize}
\item 误差分析
\end{itemize}
容量瓶中残留一定水 $\to$ 不影响；转移时有溶液溅出 $\to$ 偏小；摇匀后发现液面低于刻度线，又继续加水 $\to$ 偏小；仰视 / 俯视 $\to$ 偏小 / 偏大；砝码生锈 $\to$ 偏大；转移后不洗涤 $\to$ 偏小\par

\section{分离与提纯（物理方法）}
\subsection{过滤（液体与不溶固体 / 不溶固体与可溶固体）}
“过滤”步骤前多为“结晶”，“过滤”后多为“洗涤”、“干燥”（“蒸发”）\par

洗涤：洗涤剂上不超过滤纸边缘，下必须没过沉淀，一般洗涤 $3$ 次\par
\subsection{蒸发}
容器：蒸发皿\par

结晶方式：\par
\begin{itemize}
\item 蒸发结晶：溶解度变化小
\end{itemize}

\begin{itemize}
\item 蒸发浓缩、降温结晶：溶解度变化大
\end{itemize}

结晶后的步骤多为过滤、洗涤、干燥\par
\subsection{渗析}
用具：半透膜\par

适用范围：胶体（$d\in [1,100]$nm）\par
\subsection{分液萃取（$2$ 种互不相溶的溶液）}
目标：富集（过程为萃取）\par

操作：分层后使用分液漏斗，下层从下部流出后停住，上层由上部倒出\par

萃取剂的要求：溶解能力更强、与原溶液可分层（eg. 提取 $H_2O$ 中的 $I_2$ 可以用 $CCl_4$ 或苯作为萃取剂）\par
\begin{notebox}
补充：极性分子的正、负电荷重心不重合，离子化合物可看作极性分子\par
\end{notebox}
\subsection{蒸馏、分馏}
石油多用分馏（沸点相差 $30$℃ 以上）\par
\section{物质的鉴别}
\subsection{$Cl^-$}
加 $HNO_3$ 和 $AgNO_3$（顺序不影响），产生白色沉淀（$Ag^++Cl^-=\!=\!=AgCl\downarrow$）\par
\subsection{$SO_4^{2-}$}
先加 $HCl$（目的是除去 $CO_3^{2-}$ 和 $SO_3^{2-}$），无明显现象；再加 $BaCl_2$，产生白色沉淀（$Ba^{2+}+SO_4^{2-}=\!=\!=BaSO_4\downarrow$）\par
\subsection{$NH_4^{+}$}
先加浓 $NaOH$ 溶液，再加热，生成使红色石蕊试纸变蓝（这也是 $NH_3$ 的性质）的气体（$NH_4^++OH^-=\!=\!=NH_3\uparrow+H_2O$）\par
\begin{notebox}
扩展：$PH_3$ 也可使红色石蕊试纸变蓝\par
\end{notebox}
\subsection{$Br^-$、$I^-$}
加 $HNO_3$ 和 $AgNO_3$（顺序不影响），产生不溶于 $HNO_3$ 的（浅）黄色沉淀（$AgBr$ 是浅黄色，$AgI$ 是黄色）（$Ag^++Br^-=\!=\!=AgBr\downarrow$，\\[0.25em]
\hspace*{2em}$Ag^++I^-=\!=\!=AgI\downarrow$）\par
\subsection{$CO_3^{2-}$、$HCO_3^{-}$}
加 $HCl$，产生无色无味的气体，将气体通入澄清石灰水，澄清石灰水变浑浊（$2H^++CO_3^{2-}=\!=\!=CO_2\uparrow+H_2O$，\\[0.25em]
\hspace*{2em}$H^++HCO_3^{-}=\!=\!=CO_2\uparrow+H_2O$）\par
\subsection{$Na^+$、$K^+$}
进行焰色试验\par

$K^+$ 或 $K$ 单质：透过蓝色钴玻璃观察，火焰呈紫色\par

$Na^+$ 或 $Na$ 单质：火焰呈黄色\par
\begin{notebox}
扩展：焰色反应的实质是发射和吸收光谱\par
注：蘸取溶液进行焰色试验时，可用铁丝或铂丝（因为 $Fe$ 和 $Pt$ 没有焰色反应），铂丝更贵但更常用，因为使用 $HCl$ 清理反应后的铁丝会有 $2H^++Fe=\!=\!=Fe^{2+}+H_2\uparrow$，而 $Pt$ 则不会\par
\end{notebox}
\section{原子结构}
\subsection{道尔顿原子模型}
物质由原子构成，原子不能被创造与毁灭，原子在化学变化中不能再分割\par
\subsection{汤姆生原子模型}
汤姆生发现原子中存在电子，提出“葡萄干面包”原子模型（又叫枣糕模型）\par
\subsection{卢瑟福原子模型}
卢瑟福进行 $\alpha$ 粒子轰击金箔实验，发现原子核，提出“有核模型”\par
\subsection{玻尔原子模型}
玻尔发现 $H$ 原子光谱不连续，引入量子论观点，提出“行星模型”\par
\subsection{原子构成}
原子由原子核和电子（带负电，相对质量 $\dfrac{1}{1836}$）构成，原子核由质子（带正电，相对质量 $1.007$）和中子（不带电，相对质量 $1.008$）组成\par

$^A_ZX$：$Z$ 为质子数，$A$ 为质量数，$N=A-Z$ 为中子数，$Z$ 相同，$A$ 不同（等价于 $N$ 不同）的核素互称为同位素\par
\begin{notebox}
扩展：相对原子质量并不等于 $1.007Z+1.008N+\dfrac{1}{1836}n(e^-)$，因为核聚变会产生能量亏损，导致质量变化\par
\end{notebox}
\subsection{原子半径}
主体趋势：从左到右半径减小，从上到下半径增大\par

原理：电子屏蔽效应减弱，质子对电子吸引增强（eg. 吸引能力 $S^{2-}<Cl^-<K^+<Ca^{2+}$，电子数相同，质子数越少，吸引电子能力越弱，则半径越大）\par
\chapter{海水中的几种元素}
\section{氯}
\subsection{性质}
物理性质：黄绿色、有刺激性气味的有毒气体，密度比空气大，易液化，可溶于水\par

化学性质：有强氧化性、弱还原性\par
\subsection{常见反应}
\subsubsection{$Cl_2$ 可与除了 $Au$、$Pt$ 外的金属反应，此时 $Cl_2$ 作氧化剂：}

$2Na+Cl_2=\!=\!=2NaCl$，生成白烟\par

$2Fe+3Cl_2=\!=\!=2FeCl_3$，生成棕红色烟，且无论 $Fe$ 和 $Cl_2$ 的量有多少，生成的一定是三价铁\par

$Cu+Cl_2=\!=\!=CuCl_2$，生成棕黄色烟\par
\subsubsection{$Cl_2$可与非金属单质反应，此时 $Cl_2$ 作氧化剂：}

$H_2+Cl_2=\!=\!=2HCl$，安静地燃烧，苍白色火焰，有白雾生成\par

$2P+3Cl_2=\!=\!=2PCl_3(l)$，（$P$ 过量）\par

$2P+5Cl_2=\!=\!=2PCl_5(s)$，（$Cl_2$ 过量）\par

因此 $P$ 在 $Cl_2$ 中燃烧的现象是“有烟雾生成”，因为既有固体，也有液体生成\par
\begin{notebox}
扩展：大多数由分子构成的结构近似的物质的熔沸点与其相对分子质量正相关\par
\end{notebox}
\subsubsection{$Cl_2$ 和化合物的反应：}
$Cl_2+SO_2+2H_2O=\!=\!=2HCl+H_2SO_4$\par

$Cl_2+H_2S=\!=\!=S\downarrow+2HCl$\par

$Cl_2+2FeCl_2=\!=\!=2FeCl_3$\par

$Cl_2+H_2O\rightleftharpoons HCl+HClO$\par

$Cl_2+2NaOH=\!=\!=NaCl+NaClO+H_2O$\par

$2Cl_2+2Ca(OH)_2=\!=\!=CaCl_2+Ca(ClO)_2+2H_2O$，其中 $Ca(ClO)_2$ 是漂白粉的主要有效成分\par
\begin{notebox}
扩展：实际生产漂白粉时 $Ca(OH)_2$ 要过量，保护 $Ca(ClO)_2$ 不与 $CO_2$ 接触，以免发生 $Ca(ClO)_2+CO_2+H_2O=\!=\!=CaCO_3\downarrow +2HCl+O_2\uparrow$，这个反应也表明了 $H_2CO_3$ 的酸性强于 $HClO$\par
扩展：$Cl_2+2Na_2CO_3+H_2O=\!=\!=NaCl+NaClO+2NaHCO_3$\par
\end{notebox}
\subsection{氯水}
成分：三分子（$H_2O$、$Cl_2$、$HClO$），四离子（$H^+$、$OH^-$、$Cl^-$、$ClO^-$）\par

性质：\par
\begin{itemize}
\item 具有 $Cl_2$、$H^+$、$Cl^-$、$ClO^-$ 的性质
\end{itemize}

\begin{itemize}
\item 具有 $HClO$ 的性质：弱酸性、不稳定性、强氧化性（$HClO$ 的氧化性强于 $Cl_2$，可以杀菌、消毒、漂白）、受热或见光会分解（$2HClO=\!=\!=2HCl+O_2\uparrow$，因此保存 $HClO$ 或氯水需要棕色细口瓶，且一般放在冰箱等遮光阴凉处），因此氯水久置后酸性变强，最终变为盐酸
\end{itemize}
\begin{notebox}
扩展：“漂白”指有机色素褪色，不一定是化学反应\par
\end{notebox}
\subsection{$Cl_2$ 的制法}
\subsubsection{工业制法：}
$2NaCl+2H_2O\stackrel{\mbox{通电}}{=\!=\!=}2NaOH+H_2\uparrow+Cl_2\uparrow$\par

阳极：$2Cl^--2e^-=\!=\!=Cl_2\uparrow$\par

阴极：$2H_2O+2e^-=\!=\!=H_2\uparrow+2OH^-$\par
\subsubsection{实验室制法：}
$MnO_2+4HCl\stackrel{\Delta}{=\!=\!=}MnCl_2+2H_2O+Cl_2\uparrow$，使用饱和食盐水去除挥发的 $HCl$\par

$2KMnO_4+16HCl=\!=\!=2KCl+2MnCl_2+5Cl_2\uparrow+8H_2O$\par

$KClO_3+6HCl=\!=\!=KCl+3Cl_2+3H_2O$\par

其中，只有第一个反应需要加热，且使用 $NaOH$ 溶液吸收 $Cl_2$ 尾气\par
\section{卤素（$F$、$Cl$、$Br$、$I$）}
\subsection{物理性质}
\begin{table}[H]
\centering
\small
\renewcommand{\arraystretch}{1.12}
\begin{tabular}{c c c c}
\toprule
$F_2$ & $Cl_2$ & $Br_2$ & $I_2$ \\
\midrule
淡黄色 & 黄绿色 & 红棕色 & 紫黑色 \\
气体 & 气体 & 液体 & 固体 \\
\bottomrule
\end{tabular}
\end{table}

\subsection{化学性质}
有强氧化性，氧化能力：$F_2>Cl_2>Br_2(>Fe^{3+})>I_2$，有毒性\par
\subsection{常见反应}
\subsubsection{卤素间置换：}
$Cl_2+2Br^-=\!=\!=Br_2+2Cl^-$，$Br_2+2I^-=\!=\!=I_2+2Br^-$\par

但是 $F_2$ 不能将其它卤素离子从水溶液中置换出来，因为 $2F_2+2H_2O=\!=\!=4HF+O_2$，反应剧烈\par
\subsubsection{与 $H_2$ 反应：}
$H_2+X_2=\!=\!=2HX$，其中 $X$ 代表卤素\par
\begin{table}[H]
\centering
\scriptsize
\renewcommand{\arraystretch}{1.12}
\resizebox{0.82\linewidth}{!}{%
\begin{tabular}{c c c c}
\toprule
$F_2+H_2$ & $Cl_2+H_2$ & $Br_2+H_2$ & $I_2+H_2$ \\
\midrule
无条件爆炸 & 强光照或加热才反应 & 加热且催化才反应 & 加热且催化才反应 \\
\bottomrule
\end{tabular}%
}
\end{table}
\subsubsection{与 $H_2O$ 反应：}
氧化 $H_2O$：$2X_2+2H_2O=\!=\!=4HX+O_2$\par

歧化 $H_2O$：$X_2+H_2O=\!=\!=HX+HXO$\par
\subsubsection{与金属、非金属反应：}
大多数卤素可与非金属直接反应，所有卤素都可以与金属反应，但剧烈程度不同\par
\subsection{常见化合物}
\begin{itemize}
\item $AgF$（可溶）和 $AgCl$、$AgBr$、$AgI$（都不溶于水和硝酸，且溶解度依次递减）
\end{itemize}
\begin{notebox}
扩展：不溶于水和硝酸的沉淀还有 $BaSO_4$\par
\end{notebox}
\begin{itemize}
\item 氢卤酸
\end{itemize}
$HX$（其中 $X$ 为卤素），如 $HF$（弱酸，有强腐蚀性，有反应 $SiO_2+4HF=\!=\!=SiF_4+2H_2O$）；$HCl$、$HBr$、$HI$ 都是强酸，酸性依次递减，还原性依次增加\par

\subsection{制取方法}
$F_2$：$2HF\stackrel{\mbox{通电}}{=\!=\!=}H_2\uparrow+F_2\uparrow$\par

$Cl_2$：$2NaCl$（熔融）$\stackrel{\mbox{通电}}{=\!=\!=}2Na+Cl_2$\par

$Br_2$、$I_2$：用 $Cl_2$ 氧化 $Br^-$ 或 $I^-$\par
\section{海洋资源利用}
\subsection{$H_2O$}
实验室提取方法：蒸馏\par

仪器：蒸馏烧瓶、冷凝管（分直形、环形、球形，其中只有直形冷凝管适用于蒸馏，且水下进上出）、沸石\par
\subsection{$NaCl$}
实验室提取方法：蒸发\par

利用方法：\par
\subsubsection{制 $NaHCO_3$（侯氏制碱法）：}
$NH_3+H_2O=\!=\!=NH_3\cdot H_2O$\par

$CO_2+NH_3\cdot H_2O=\!=\!=NH_4HCO_3$\par

$NH_4HCO_3+NaCl=\!=\!=NH_4Cl+NaHCO_3\downarrow$\par

$2NaHCO_3\stackrel{\Delta}{=\!=\!=}Na_2CO_3+H_2O+CO_2\uparrow$\par

总反应：$NH_3+H_2O+CO_2+NaCl=\!=\!=NaHCO_3+NH_4Cl$\par
\begin{notebox}
注意：先加 $NH_3$ 再加 $CO_2$，可以提高 $CO_2$ 溶解度\par
\end{notebox}
\subsubsection{制 $Na$ 与 $Cl_2$：}
$2NaCl$（熔融）$\stackrel{\mbox{通电}}{=\!=\!=}2Na+Cl_2\uparrow$\par

\subsubsection{制消毒剂：}
$2NaCl+2H_2O\stackrel{\mbox{通电}}{=\!=\!=}2NaOH+H_2\uparrow+Cl_2\uparrow$\par

$Cl_2+2NaOH=\!=\!=NaCl+NaClO+H_2O$\par
\subsection{$MgCl_2$}
提纯：思路是拿出目标物再转化（先碱后酸）\par

$MgCl_2+Ca(OH)_2$（浊液）$=\!=\!=Mg(OH)_2\downarrow+2CaCl_2$\par

$Mg(OH)_2+2HCl=\!=\!=MgCl_2+H_2O$\par
\subsection{$I_2$}
提纯：从海带中提取\par

先把海带烧成灰（使用仪器：坩埚），加水溶解、过滤，再利用反应 $Cl_2+2I^-=\!=\!=I_2+2Cl^-$，最后蒸馏（从 $I_2$ 的 $CCl_4$ 溶液中蒸馏出）\par
\section{钠}
\subsection{性质}
物理性质：银白色固体，导电，质地软，密度比水小（$0.97$g/cm$^3$），熔点低\par

化学性质：强还原性，能与所有非金属反应，火焰黄色\par
\subsection{常见反应}
\subsubsection{与 $H_2O$ 反应：}
$2Na+2H_2O=\!=\!=2NaOH+H_2\uparrow$\par

（浮游熔响红）钠熔成小球，浮在水面，四处游动，发出“嘶嘶”的响声，（若水中有酚酞，则溶液会变红）\par
\subsubsection{与 $CuSO_4$ 溶液反应：}
$2Na+2H_2O=\!=\!=2NaOH+H_2\uparrow$\par

$2NaOH+CuSO_4=\!=\!=Na_2SO_4+Cu(OH)_2\downarrow$\par

$Cu(OH)_2\stackrel{\Delta}{=\!=\!=}CuO+H_2O$（前面的反应放热）\par

综上，反应生成蓝色（即 $Cu(OH)_2$）和黑色（即 $CuO$）沉淀\par

总反应：$2Na+CuSO_4+2H_2O=\!=\!=Na_2SO_4+Cu(OH)_2\downarrow+H_2\uparrow$\par

离子方程式：$2Na+Cu^{2+}+2H_2O=\!=\!=2Na^++Cu(OH)_2\downarrow+H_2\uparrow$\par
\subsubsection{与 $O_2$ 反应：}
$4Na+O_2=\!=\!=2Na_2O$\par

$2Na+O_2\stackrel{\mbox{点燃}}{=\!=\!=}Na_2O_2$，其中，$Na_2O_2$ 是黄色的，$O_2^{2-}$ 是过氧根离子团，现象是有黄色火焰\par
\begin{notebox}
注：$1$mol 的 $Na_2O$ 和 $Na_2O_2$ 都含 $3$mol 离子，且正负离子比为 $2:1$\par
\end{notebox}
在空气中的后续反应：最终都变为 $Na_2CO_3$\par

$Na_2O+H_2O=\!=\!=2NaOH$，$2NaOH+CO_2=\!=\!=Na_2CO_3+H_2O$\par

$2Na_2O_2+2H_2O=\!=\!=4NaOH+O_2\uparrow$，\\[0.25em]
\hspace*{2em}$2NaOH+CO_2=\!=\!=Na_2CO_3+H_2O$\par

$Na_2O_2$ 有强氧化性，且可用于呼吸面罩制氧气（吸收 $CO_2$，生成 $O_2$）\par
\subsubsection{与 $Cl_2$ 反应：}
$2Na+Cl_2=\!=\!=2NaCl$，无论是否点燃都可以反应，现象是发出白烟\par
\subsubsection{与 $TiCl_4$ 反应：}
$4Na+TiCl_4$（熔融）$\stackrel{\mbox{高温}}{=\!=\!=}Ti+4NaCl$\par

\begin{notebox}
注：$Na$ 在盐溶液中会先和 $H_2O$ 反应，不会置换，但是 $Na$ 可以与盐的熔融态置换\par
\end{notebox}
\subsection{用途}
冶炼金属、使用钠钾合金（液体、比热容大于水）作核电站导热剂、钠灯（透雾、穿透力强）\par
\subsection{保存方法}
密封保存（少量的 $Na$ 可保存于煤油中）\par
\begin{notebox}
注：$K$ 和 $Na$ 的密度都比煤油大，所以可以保存在煤油中，但是 $Li$ 的密度比煤油小，所以封在石蜡中保存\par
\end{notebox}
\subsection{制取}
$2NaCl\stackrel{\mbox{通电}}{=\!=\!=}2Na+Cl_2\uparrow$\par

阳极反应：$2Cl^--2e^-=\!=\!=Cl_2\uparrow$\par

阴极反应：$2Na^++2e^-=\!=\!=2Na$\par
\section{钠的化合物}
几乎所有钠的化合物都是电解质和离子化合物\par

\subsection{氧化物（$Na_2O$、$Na_2O_2$）}
$Na_2O$：碱性氧化物（$Na_2O+H_2O=\!=\!=2NaOH$）\par

$Na_2O_2$：其它氧化物，具有氧化性，可以杀菌消毒和漂白\par
\begin{notebox}
注：漂白是使有机色素变为无色，pH 试纸含有机色素，故有漂白性的物质不能用 pH 试纸测\par
\end{notebox}
$2Na_2O_2+2H_2O=\!=\!=4NaOH+O_2\uparrow$\par

$2Na_2O_2+2CO_2=\!=\!=2Na_2CO_3+O_2$\par

以上两个反应说明 $Na_2O_2$ 可作潜水供氧剂\par

$Na_2O_2$ 也可以与酸反应：$2Na_2O_2+4HCl=\!=\!=4NaCl+2H_2O+O_2\uparrow$\par

\begin{notebox}
扩展：$Na_2O_2+2H_2O=\!=\!=H_2O_2+2NaOH$，证明产生 $H_2O_2$ 的方法：加过氧化氢酶\par
\end{notebox}
\subsection{苏打}
\begin{table}[H]
\centering
\scriptsize
\renewcommand{\arraystretch}{1.14}
\resizebox{\linewidth}{!}{%
\begin{tabular}{c c c c c c}
\toprule
 & 俗名 & 溶液 pH & 溶解度 & 受热 & 与酸反应 \\
\midrule
$NaHCO_3$ & 小苏打 & pH 约为 $8$ & 可溶（约 $8.4$g） & 易分解 & 快 \\
$Na_2CO_3$ & 纯碱、苏打 & pH 约为 $12$ & 易溶（约 $24$g） & 稳定 & 慢 \\
$Na_2S_2O_3$ & 大苏打、海波 & / & / & / & / \\
\bottomrule
\end{tabular}%
}
\end{table}

与 $CaCl_2$ 反应：\par

$2NaHCO_3+CaCl_2=\!=\!=2NaCl+Ca(HCO_3)_2$，其中 $Ca(HCO_3)_2$ 微溶，故加入少量 $CaCl_2$ 时无明显现象\par

$Na_2CO_3+CaCl_2=\!=\!=2NaCl+CaCO_3\downarrow$，其中 $CaCO_3$ 难溶，会立即产生白色沉淀\par

转化关系：\par

$2NaHCO_3\stackrel{\Delta}{=\!=\!=}Na_2CO_3+H_2O+CO_2\uparrow$\par

$Na_2CO_3+H_2O+CO_2=\!=\!=2NaHCO_3$，若 $Na_2CO_3$ 是饱和的，则一定有 $NaHCO_3$ 析出\par

鉴别 $Na_2CO_3$ 和 $NaHCO_3$ 的方法：\par

1. 看标签\par

2. 看瓶塞（$Na_2CO_3$ 溶液强碱性，用橡胶塞，$NaHCO_3$ 用玻璃塞）\par

3. 加热，观察是否分解\par

4. 测饱和溶液 pH\par

5. 与盐酸反应，观察速率\par

6. 溶解一定量，看是否全溶\par

7. 缓慢加 $CaCl_2$，看现象\par

8. 等量固体加盐酸，看产生 $CO_2$ 的量\par

\section{金属的冶炼方法（补充）}
\subsection{活泼金属}
$K$、$Ca$、$Na$、$Mg$、$Al$ 通常用电解法，电解时常会加一些其他物质来降低熔化温度\par

\begin{table}[H]
\centering
\scriptsize
\renewcommand{\arraystretch}{1.14}
\resizebox{\linewidth}{!}{%
\begin{tabular}{c c c c c c}
\toprule
产物 & $K$ & $Ca$ & $Na$ & $Mg$ & $Al$ \\
\midrule
原料（熔融） & $KCl$ & $CaCl_2$ & $NaCl$ & $MgCl_2$ & $Al_2O_3$ \\
\bottomrule
\end{tabular}%
}
\end{table}

\subsection{$Zn$、$Fe$、$Sn$、$Pb$、$Cu$}
通常使用还原法，常用 $C$ 或 $CO$ 作还原剂\par
\begin{notebox}
扩展：特殊还原剂 $4Na+TiCl_4=\!=\!=Ti+4NaCl$\par
\end{notebox}
另外， $Al$ 几乎可以置换出所有的金属单质（甚至可以置换出 $K$、$Ca$、$Na$、$Mg$），且反应放热（铝热反应）\par
\subsection{$Ag$、$Hg$}
加强热，使其受热分解\par

含 $Ag$、$Hg$ 的化合物（氟化物除外）都可以加强热以得到其单质\par
\subsection{特殊方法}
$Na+KCl\rightleftharpoons K\uparrow +NaCl$，气态的 $K$ 飘走，使反应向右进行，且 $Na$ 熔沸点比 $K$ 高\par

\section{化学反应的本质（补充）}
$\Delta H$：焓（生成物与反应物能量之差）\par

$\Delta S$：熵（混乱度），物质由固体变为液体再变为气体熵增，气体从无到有或增多为熵增，反之亦然\par

结论：当 $\Delta H-T \Delta S<0$ 时，化学反应可行（$T$ 是开氏温度）\par

推论：\par

\begin{itemize}
\item $\Delta H<0$，$\Delta S>0$ 时，反应不可逆，任意条件可行
\end{itemize}

\begin{itemize}
\item $\Delta H>0$，$\Delta S<0$ 时，反应不可逆，永不可行
\end{itemize}

\begin{itemize}
\item $\Delta H<0$，$\Delta S<0$ 时，反应可逆，低温时可行
\end{itemize}

\begin{itemize}
\item $\Delta H>0$，$\Delta S>0$ 时，反应可逆，高温时可行
\end{itemize}

技巧：\par

1. 吸热 $\Delta H>0$，放热 $\Delta H<0$\par

2. 通过物态变化判断 $\Delta S$\par

3. $\Delta H$、$\Delta S$ 同号则反应可逆，异号则不可逆\par
\chapter{硫与环境保护}
\section{硫单质}
\subsection{性质}
物理性质：非金属，黄色晶体（粉末状沉淀为乳白色），分子晶体，块状，不溶于水，微溶于酒精，易溶于 $CS_2$\par

化学性质：有弱氧化性、弱还原性\par
\subsection{反应}
\subsubsection{氧化性：}
除 $Au$、$Pt$ 外的所有金属都可以和 $S$ 反应\par

$Hg+S=\!=\!=HgS$，俗名丹砂、朱砂\par

$Fe+S\stackrel{\Delta}{=\!=\!=}FeS$，\\[0.25em]
\hspace*{2em}$2Cu+S\stackrel{\Delta}{=\!=\!=}Cu_2S$，金属都不是最高价\par

$H_2+S\stackrel{\Delta}{=\!=\!=}H_2S$，无色，有臭鸡蛋味，有毒\par
\subsubsection{还原性：}
$S+O_2\stackrel{\mbox{点燃}}{=\!=\!=}SO_2$\par

$S+3F_2=\!=\!=SF_6$\par

$2S$（多）$+Cl_2\stackrel{\Delta}{=\!=\!=}S_2Cl_2$；$S$（少）$+Cl_2\stackrel{\Delta}{=\!=\!=}SCl_2$\par

$S+2H_2SO_4$（浓）$=\!=\!=3SO_2\uparrow+2H_2O$\par

$S+6HNO_3$（浓）$=\!=\!=H_2SO_4+6NO_2\uparrow+2H_2O$\par

$S+2HNO_3=\!=\!=H_2SO_4+2NO\uparrow$\par
\subsubsection{其他反应：}
$3S+6NaOH\stackrel{\Delta}{=\!=\!=}2Na_2S+Na_2SO_3+3H_2O$\par

$Na_2S+(x-1)S\rightleftharpoons Na_2S_x$（多硫化钠，$x_{max}=5$）\par

$Na_2SO_3+S\stackrel{\Delta}{\rightleftharpoons}Na_2S_2O_3$\par

$SO_2+Cl_2+2H_2O=\!=\!=2HCl+H_2SO_4$\par
\subsection{用途}
\begin{itemize}
\item 作为化工原料（eg. 制硫酸）
\end{itemize}

\begin{itemize}
\item 做药膏或火药（$S+2KNO_3+3C=\!=\!=3CO_2+K_2S+N_2\uparrow$）
\end{itemize}

\section{$H_2S$}
\subsection{性质}
物理性质：无色、臭鸡蛋味气体，密度略大于空气，能溶于水，有剧毒\par

化学性质：有酸性（略弱于 $H_2CO_3$，强于 $HClO$），有弱氧化性和强还原性，不稳定\par
\subsection{常见反应}
$H_2S+CuSO_4=\!=\!=CuS\downarrow+H_2SO_4$\par

$H_2S+Cu(OH)_2=\!=\!=CuS\downarrow+2H_2O$\par

$H_2S+X_2\stackrel{\Delta}{=\!=\!=}S+2HX$，$X$ 为卤素\par

$2H_2S+SO_2=\!=\!=3S\downarrow+2H_2O$，归中反应\par

$H_2S+H_2SO_4$（浓）$=\!=\!=S\downarrow+SO_2+2H_2O$\par

$2H_2S+O_2=\!=\!=2S\downarrow+2H_2O$，在水中时\par

$2H_2S+3O_2\stackrel{\Delta}{=\!=\!=}2SO_2+2H_2O$，在空气中且有足量 $O_2$ 时\par

$2H_2S+O_2\stackrel{\Delta}{=\!=\!=}2S+2H_2O$，在空气中且有少量 $O_2$ 时\par

$2Fe^{3+}+H_2S=\!=\!=S\downarrow+2Fe^{2+}+2H^+$\par
\subsection{制取方法}
$FeS+H_2SO_4=\!=\!=H_2S\uparrow+FeSO_4$\par

制取装置：启普发生器（可以和制取 $CO_2$ 使用同一装置）\par

收集方法：向上排空气法\par
\section{$SO_3$}
\subsection{性质}
物理性质：标况（$0$℃）为固体，常温（$25$℃）为液体，沸点 $44.5$℃，无色\par

化学性质：有氧化性（高温下非常强）和还原性，是酸性氧化物，能与碱、碱性氧化物和水反应\par
\subsection{常见反应}
$SO_3+2OH^-=\!=\!=SO_4^{2-}+H_2O$，$OH^-$ 多时\par

$SO_3+OH^-=\!=\!=HSO_4^-$，$SO_3$ 多时\par

$SO_3+H_2O=\!=\!=H_2SO_4$\par
\begin{notebox}
扩展：$P_2O_5$ 是唯一强于 $H_2SO_4$ 的吸水剂，也是目前最强的吸水剂\par
\end{notebox}
$2SO_3\xrightleftharpoons[\mbox{\mbox{催化剂}}]{\Delta}2SO_2+O_2$\par
\section{$SO_2$}
\subsection{性质}
物理性质：无色、有毒、有刺激性气味的气体，易溶于水（按 $1:40$ 的比例溶解）\par

化学性质：酸性氧化物，有漂白性（能使品红溶液褪色，加热后恢复）\par
\subsection{常见反应}
\subsubsection{氧化性：}
$SO_2+2CO\stackrel{\Delta}{=\!=\!=}S+2CO_2$\par
\subsubsection{还原性：}
$2Fe^{3+}+SO_2+2H_2O=\!=\!=2Fe^{2+}+SO_4^{2-}+4H^+$\par

$SO_2+Cl_2+2H_2O=\!=\!=H_2SO_4+2HCl$，这个反应说明了：\par

\begin{itemize}
\item $SO_2$ 和 $Cl_2$ 都有漂白性，但不能同时使用，因为会反应
\end{itemize}

\begin{itemize}
\item $SO_2$ 使氯水、溴水、高锰酸钾溶液褪色不体现其漂白性
\end{itemize}

\begin{notebox}
扩展：$SO_2+Cl_2\stackrel{\Delta}{=\!=\!=}SO_2Cl_2$，生成物中的 $S$ 为 $+6$ 价，$Cl$ 为 $-1$ 价，这说明 $SO_2$ 与 $Cl_2$ 在无水时也可以反应\par
\end{notebox}
\subsubsection{制硫酸：}
$SO_2\stackrel{H_2O}{\longrightarrow}H_2SO_3\stackrel{O_2}{\longrightarrow}H_2SO_4$，可自然发生（酸雨的形成），浓度低\par

$SO_2\stackrel{O_2}{\longrightarrow}SO_3\stackrel{H_2O}{\longrightarrow}H_2SO_4$，需要加热，浓度更高，是工业制法\par
\subsection{用途}
\begin{itemize}
\item 杀菌消毒（eg. 给葡萄酒作消毒剂）、抗氧化
\end{itemize}

\begin{itemize}
\item 作为化工原料（eg. 制 $H_2SO_4$）
\end{itemize}
\subsection{制法}
\subsubsection{工业制法：}
外国：$S+O_2=\!=\!=SO_2$\par

中国：$3FeS_2+8O_2=\!=\!=Fe_3O_4+6SO_2$，其中 $FeS_2$ 要先粉碎\par
\subsubsection{实验室制法：}
$Na_2SO_3+H_2SO_4$（中等浓度）$\stackrel{\Delta}{=\!=\!=}Na_2SO_4+H_2O+SO_2\uparrow$\par
\subsubsection{尾气吸收：}
用 $Na_2SO_3$（工业）或 $NaOH$（实验室）\par
\section{$H_2SO_4$}
$1$mol$\cdot$L$^{-1}$ 的硫酸的氢离子浓度约为 $1.03$mol$\cdot$L$^{-1}$\par
\subsection{性质}
\subsubsection{作为稀酸的通性：}
可与指示剂反应（使石蕊变红、甲基橙变红）、与活泼金属反应、与碱中和、与碱性（两性）氧化物反应、与碱式盐和其它的部分盐（eg. $CO_3^{2-}$、$HCO_3^-$）反应\par

\subsubsection{浓硫酸特性：}
\begin{itemize}
\item 吸水性（物理性质，可作为气体干燥剂，但不能干燥除了 $SO_2$ 的碱性和还原性气体）
\end{itemize}

\begin{itemize}
\item 脱水性（化学性质）
\end{itemize}

（eg.$C_{12}H_{22}O_{11}\stackrel{\mbox{浓} H_2SO_4}{\longrightarrow}12C+11H_2O$）\par

\begin{itemize}
\item 强氧化性
\end{itemize}

$S$ 从 $+6$ 价降为 $+4$ 价，一般从 $SO_4^{2-}$ 变为 $SO_2$，它可以与除了 $Au$ 和 $Pt$ 之外的所有金属反应(eg. $Cu+2H_2SO_4$（浓）$\stackrel{\Delta}{=\!=\!=}CuSO_4+SO_2\uparrow+2H_2O$\par

另外，它还可以与部分非金属反应（eg. $C+2H_2SO_4$（浓）$\stackrel{\Delta}{=\!=\!=}CO_2\uparrow+2SO_2\uparrow+2H_2O$，\\[0.25em]
\hspace*{2em}$S+2H_2SO_4=\!=\!=3SO_2\uparrow+2H_2O$）\par

\begin{itemize}
\item 钝化（强氧化性的体现）
\end{itemize}

$Fe$、$Al$ 放入浓硫酸后，会形成致密的氧化膜（$Fe_3O_4$、$Al_2O_3$），就不会再继续反应\par

\begin{notebox}
扩展：黑面包实验\par
$C_{12}H_{22}O_{11}\stackrel{\mbox{浓} H_2SO_4}{\longrightarrow}12C+11H_2O$，$C+2H_2SO_4$（浓）$\stackrel{\Delta}{=\!=\!=}CO_2\uparrow+2SO_2\uparrow+2H_2O$\par
现象：先变黄后变黑，再膨胀\par
若想增大反应速率，可以加入几滴水来提高温度\par
\end{notebox}
\subsection{制取}
1. $2SO_2+O_2\xrightleftharpoons[V_2O_5]{\Delta}2SO_3$，\\[0.25em]
\hspace*{2em}$SO_3+H_2O=\!=\!=H_2SO_4$（这一步需要用浓硫酸吸收）\par

2. $SO_2+H_2O=\!=\!=H_2SO_3$，$2H_2SO_3+O_2=\!=\!=2H_2SO_4$\par

3. $2FeSO_4\cdot 7H_2O\stackrel{\Delta}{=\!=\!=}Fe_2O_3+SO_2\uparrow+SO_3\uparrow+14H_2O$，\\[0.25em]
\hspace*{2em}$SO_3+H_2O=\!=\!=H_2SO_4$\par

4. $CuSO_4\cdot 5H_2O\stackrel{\Delta}{=\!=\!=}CuO+SO_3\uparrow+5H_2O$，\\[0.25em]
\hspace*{2em}$SO_3+H_2O=\!=\!=H_2SO_4$\par

\begin{notebox}
注：后两种方法生成的硫酸浓度较小\par
\end{notebox}
\section{氮气}
\subsection{性质}
物理性质：无色无味的无毒气体，难溶于水\par

化学性质：通常很稳定，结构：$N\equiv N$\par
\subsection{常见反应}
\subsubsection{氧化性：}
$N$ 的最低价为 $-3$ 价\par

$6Li+N_2=\!=\!=2Li_3N$，这个反应引出了过渡态理论\par

$3Mg+N_2\stackrel{\mbox{点燃}}{=\!=\!=}Mg_3N_2$\par

$N_2+3H_2\xrightleftharpoons[\mbox{\mbox{催化剂}}]{\mbox{\mbox{高温高压}}}2NH_3$\par
\subsubsection{还原性：}
$N$ 的最高价是 $+5$ 价\par

$N_2+O_2\stackrel{\mbox{放电}}{=\!=\!=}2NO$，后续立即被 $O_2$ 氧化：$2NO+O_2=\!=\!=2NO_2$，$NO_2$ 是红棕色气体\par
\subsubsection{固氮反应：}
把 $N_2$ 转化为氮的化合物，体现氮的氧化性，常见的有豆科植物根瘤菌\par
\subsection{用途}
合成氨、制硝酸、作为化工原料、制冷剂、作为保护气\par
\section{氨气}
\subsection{性质}
物理性质：无色、有刺激性气味的气体，极易溶于水（因为是极性分子且可以与水形成分子间氢键，溶解比例 $1:700$），氨水的密度 $<1$，且浓度越大，密度越小\par

化学性质：有弱氧化性、强还原性，结构为三角锥性，键角 $107^{\circ}18'$，有碱性（鉴别氨气：用湿润的红色石蕊试纸靠近，试纸变蓝）\par
\subsection{常见反应}
\subsubsection{氧化性：}
$2K+2NH_3(l)=\!=\!=H_2\uparrow+2KNH_2$\par
\subsubsection{还原性：}
$4NH_3+5O_2\stackrel{Pt}{=\!=\!=}4NO+6H_2O$（催化氧化）\par

$NH_3+Cl_2\rightarrow N_2+HCl/NH_4Cl$\par

$4NH_3+6NO=\!=\!=5N_2+6H_2O$\par
\subsection{制法}
工业：$N_2+3H_2\xrightleftharpoons[\mbox{\mbox{催化剂}}]{\mbox{\mbox{高温高压}}}2NH_3$\par

实验室（少量）：$2NH_4Cl+Ca(OH)_2\stackrel{\Delta}{=\!=\!=}CaCl_2+2NH_3\uparrow +2H_2O$\par
\subsection{用途}
制化肥、制硝酸、工业原料\par
\section{氮的氧化物}
\subsection{种类}
\begin{table}[H]
\centering
\renewcommand{\arraystretch}{1.22}
\begin{tabularx}{\linewidth}{Y Y Y Y Y Y}
\toprule
价态 & $+1$ & $+2$ & $+3$ & $+4$ & $+5$ \\
\midrule
化学式 & $N_2O$ & $NO$、$N_2O_2$ & $N_2O_3$ & $NO_2$、$N_2O_4$ & $N_2O_5$ \\
\bottomrule
\end{tabularx}
\end{table}

其中，$N_2O_3$ 是亚硝酸酸酐，$N_2O_5$ 是硝酸酸酐，它们都是酸性氧化物\par
\subsection{$NO$}
\subsubsection{性质：}
$NO$ 是大气污染物\par

物理性质：密度略小于空气，无色气体，微溶于水\par

化学性质：有弱氧化性、强还原性、不稳定性和毒性\par
\subsubsection{常见反应：}
\begin{itemize}
\item 氧化性：
\end{itemize}

$4NH_3+6NO=\!=\!=5N_2+6H_2O$\par

$2CO+2NO=\!=\!=N_2+2CO_2$\par

\begin{itemize}
\item 还原性：
\end{itemize}

$2NO+O_2=\!=\!=2NO_2$\par

$4NO+3O_2+2H_2O=\!=\!=4HNO_3$\par
\subsubsection{危害：}
造成酸雨、臭氧层空洞和光化学烟雾\par
\subsubsection{用途：}
速效救心丸\par
\subsection{$NO_2$}
\subsubsection{性质：}
$NO_2$ 是大气污染物\par

物理性质：常温时是红棕色刺鼻气体\par

化学性质：有强氧化性、弱还原性，较稳定\par
\subsubsection{常见反应：}
\begin{itemize}
\item 氧化性（室温下强于 $Cl_2$）：
\end{itemize}

$NO+NO_2+2NaOH=\!=\!=2NaNO_2+H_2O$\par

\begin{itemize}
\item 还原性：
\end{itemize}

$3NO_2+H_2O=\!=\!=2HNO_3+NO$\par

$2NO_2+NaOH=\!=\!=NaNO_3+NaNO_2+H_2O$\par
\subsubsection{用途：}
制硝酸\par
\subsubsection{危害：}
光化学烟雾\par
\section{$HNO_3$}
\subsection{性质}
物理性质：易溶于水，常温下其水溶液无色透明\par

化学性质：有酸性、不稳定性、强氧化性\par
\subsection{保存}
棕色细口瓶，放置于阴凉处\par
\subsection{常见反应}
\subsubsection{分解：}
$4HNO_3\stackrel{\Delta /\mbox{光照}}{=\!=\!=}4NO_2\uparrow+O_2\uparrow+2H_2O$\par
\subsubsection{氧化性：}
$3Cu+8HNO_3=\!=\!=3Cu(NO_3)_2+2NO\uparrow+4H_2O$\par

$Cu+4HNO_3$（浓）$=\!=\!=Cu(NO_3)_2+2NO_2\uparrow+2H_2O$\par

$C+4HNO_3$（浓）$\stackrel{\Delta}{=\!=\!=}CO_2\uparrow+4NO_2\uparrow+2H_2O$\par
\subsubsection{钝化：}
可以钝化的金属有：$Fe$、$Al$、$Cr$、$Ni$ 等\par

$Fe$ 会变成 $Fe_2O_3$，$Al$ 会变成 $Al_2O_3$\par

\begin{notebox}
扩展：$Fe$ 与 $HNO_3$ 反应\par
$Fe$（少）$+HNO_3$（多）$\rightarrow Fe(NO_3)_3+NO\uparrow+H_2O$\par
$Fe$（多）$+HNO_3$（少）$\rightarrow Fe(NO_3)_2+NO\uparrow+H_2O$\par
$Fe$（少）$+HNO_3$（浓、多）$\stackrel{\Delta}{\rightarrow} Fe(NO_3)_3+NO_2\uparrow+H_2O$\par
$Fe$（多）$+HNO_3$（浓、少）$\stackrel{\Delta}{\rightarrow} Fe(NO_3)_2+NO_2\uparrow+H_2O$\par
反应物是浓硝酸时会钝化，所以需要加热\par
$Fe$ 过量时，反应的实质是 $2Fe^{3+}+Fe=\!=\!=3Fe^{2+}$\par
\end{notebox}
\subsection{制法}
\subsubsection{实验室制法：}
$NaNO_3+H_2SO_4$（浓）$\stackrel{\Delta}{=\!=\!=}HNO_3\uparrow+NaHSO_4$\par
\subsubsection{工业制法：}
$4NH_3+5O_2\stackrel{\Delta}{=\!=\!=}4NO+6H_2O$\par

$2NO+O_2=\!=\!=2NO_2$，$H_2O+3NO_2=\!=\!=2HNO_3+NO$（循环氧化）\par

第二步总反应：$4NO+3O_2+2H_2O=\!=\!=4HNO_3$\par
\chapter{其它重要元素}
\section{铁}
\subsection{$Fe$（单质）}
$2Fe+3Cl_2/Br_2=\!=\!=2Fe^{3+}+6Cl^-/Br^-$\par

$Fe+I_2/S=\!=\!=FeI_2/FeS$\par

$4Fe+3O_2\stackrel{\Delta}{=\!=\!=}2Fe_2O_3$，\\[0.25em]
\hspace*{2em}$3Fe+2O_2\stackrel{\mbox{点燃}}{=\!=\!=}Fe_3O_4$\par

$3Fe+4H_2O$（g）$\stackrel{\Delta}{\rightleftharpoons}Fe_3O_4+4H_2$\par
\subsection{$Fe^{2+}$}
\subsubsection{氧化性：}
$FeO+H_2\stackrel{\Delta}{=\!=\!=}Fe+H_2O$\par

$FeO+C\stackrel{\Delta}{=\!=\!=}Fe+CO\uparrow$\par

$FeO+CO\stackrel{\Delta}{=\!=\!=}Fe+CO_2$\par

$FeO+NaBH_4\rightarrow Fe$\par
\subsubsection{还原性：}
$2Fe^{2+}+H_2O_2+4OH^-=\!=\!=2Fe(OH)_3\downarrow$\par

$2Fe^{2+}+H_2O_2+2H^+=\!=\!=2Fe^{3+}+2H_2O$\par

$Fe^{2+}+Cl_2/Br_2/HNO_3\rightarrow Fe^{3+}$\par
\subsection{$Fe^{3+}$}
$2Fe^{3+}+Fe=\!=\!=3Fe^{2+}$\par

$2Fe^{3+}+Cu=\!=\!=Cu^{2+}+2Fe^{2+}$\par

$Fe^{3+}+I^-/S^{2-}\rightarrow Fe^{2+}$（$I^-\rightarrow I_2$，$S^{2-}\rightarrow S$）\par

$Fe^{3+}+SCN^-=\!=\!=Fe(SCN)_3$（血红色，可以用于检验 $Fe^{3+}$）\par

$Fe^{3+}+[Fe(CN)_6]^{4-}$ 会变蓝，可用于检验 $Fe^{3+}$\par
\subsection{检验 $Fe^{2+}$ 和 $Fe^{3+}$}
\subsubsection{证明存在 $Fe^{2+}$：}
先加 $KSCN$，无明显现象，再加氧化剂，变红色\par
\subsubsection{证明同时存在 $Fe^{2+}$ 和 $Fe^{3+}$：}
先取两份样品 A 和 B,\par

1. A 加 $KSCN$，出现血红色，则有 $Fe^{3+}$；B 加 $KMnO_4$，紫色褪去，则有 $Fe^{2+}$\par

2. A 加 $K_4[Fe(CN)_6]$，变蓝，则有 $Fe^{3+}$；B 加 $K_3[Fe(CN)_6]$，变蓝，则有 $Fe^{2+}$\par

\subsection{冶炼铁的氧化物}
$Fe_2O_3+3CO\stackrel{\Delta}{\rightleftharpoons}2Fe+3CO_2$，这个反应也揭示了可逆反应的存在\par
\section{硅}
\subsection{Si（单质）}
硅单质是共价晶体\par

\begin{itemize}
\item 常见反应：
\end{itemize}

$Si+2Cl_2\stackrel{\Delta}{=\!=\!=}SiCl_4$\par

$Si+4HF=\!=\!=SiF_4+2H_2\uparrow$\par

$H_2O+Si+2NaOH=\!=\!=Na_2SiO_3+2H_2\uparrow$\par

\begin{itemize}
\item 用途：太阳能电池、晶体管
\end{itemize}
\subsection{$SiO_2$}
$SiO_2$ 是酸性氧化物，可用于生产光纤\par

$SiO_2+4HF=\!=\!=SiF_4+2H_2O$（事实上不是）\par

$SiO_2+6HF=\!=\!=H_2SiF_6+2H_2O$（真实情况）\par

$SiO_2+4H^++6F^-=\!=\!=SiF_6^{2-}+2H_2O$（离子方程式）\par

$SiO_2+2NaOH=\!=\!=Na_2SiO_3+H_2O$\par

$SiO_2+3C\stackrel{\mbox{高温}}{=\!=\!=}SiC+2CO\uparrow$，\\[0.25em]
\hspace*{2em}$SiO_2+2C\stackrel{\mbox{高温}}{=\!=\!=}Si+2CO\uparrow$\par
\subsection{硅的提纯}
$Si+2Cl_2=\!=\!=SiCl_4$，$SiCl_4+2H_2=\!=\!=Si+4HCl$\par
\section{铜}
\subsection{Cu（单质）}
$2Cu+S\stackrel{\Delta}{=\!=\!=}Cu_2S$\par
\subsection{$Cu^{2+}$}
$2Cu^{2+}+4I^-=\!=\!=2CuI\downarrow+I_2$\par

$Cu^{2+}$ 可以氧化 $S^{2-}$，所以 $CuS$ 中各元素价态较复杂\par
\chapter{结构化学}
\section{元素周期律（上）}
\subsection{原子核外电子排布呈周期性变化}
\subsection{化合价呈周期性变化（最高正价$-$最低负价$=8$）}
\subsection{原子半径呈周期性变化}
同一周期从左到右原子半径减小，同一族从上到下原子半径增大（最小的是 $H$，$53$pm；第二小的是 $F$，$64$pm）\par

原子半径表（单位为 pm，即 $10^{-12}$m）：\par
\begin{table}[H]
\centering
\renewcommand{\arraystretch}{1.22}
\begin{tabularx}{\linewidth}{Y Y Y Y Y Y Y}
\toprule
$Li$ & $Be$ & $B$ & $C$ & $N$ & $O$ & $F$ \\
\midrule
$157$ & $112$ & $88$ & $77$ & $74$ & $66$ & $64$ \\
$Na$ & $Mg$ & $Al$ & $Si$ & $P$ & $S$ & $Cl$ \\
$191$ & $160$ & $143$ & $118$ & $110$ & $104$ & $79$ \\
\bottomrule
\end{tabularx}
\end{table}

常见的离子半径比较：\par

$O^{2-}>Na^+>Mg^{2+}>Al^{3+}$\par

$Fe^{2+}>Fe^{3+}$，$O^{2-}>O$，$Na>Na^+$\par
\subsection{金属性（原子失电子能力）}
\subsubsection{强弱鉴别：}
1. （中等活泼元素）与非氧化性酸反应，越剧烈，金属性越强\par

2. （极活泼元素）与水反应，越剧烈，金属性越强\par

3. （不活泼元素）发生置换反应，反应物单质强于生成物单质\par

4. 与最高价氧化物对应的碱的碱性正相关\par
\subsubsection{变化规律：}
同周期从左到右，金属性依次减弱；同族从上到下，金属性依次增强\par
\subsection{非金属性（原子得电子能力）}
\subsubsection{强弱鉴别：}
1.  与 $H_2$ 反应，越剧烈，非金属性越强\par

2.  置换反应中，生成物单质强于反应物单质\par

3.  与氢化物稳定性正相关\par

4.  与最高价含氧酸的酸性正相关\par
\subsubsection{变化规律：}
同周期从左到右，非金属性依次增强；同族从上到下，非金属性依次减弱\par
\section{原子核外电子排布}
\subsection{轨道}
\begin{itemize}
\item s：球形，有 $1$ 个轨道，可容纳 $2$ 个电子
\end{itemize}

\begin{itemize}
\item p：纺锤形，有 $3$ 个轨道，可容纳 $6$ 个电子
\end{itemize}

\begin{itemize}
\item d：有 $5$ 个轨道，可容纳 $10$ 个电子
\end{itemize}

\begin{itemize}
\item f：有 $7$ 个轨道，可容纳 $14$ 个电子
\end{itemize}
\subsection{电子在轨道上的排布}
\begingroup
\small
\setlength{\tabcolsep}{3pt}
\renewcommand{\arraystretch}{1.12}
\begin{longtable}{>{\centering\arraybackslash}p{0.16\linewidth} >{\centering\arraybackslash}p{0.14\linewidth} >{\centering\arraybackslash}p{0.43\linewidth} >{\centering\arraybackslash}p{0.16\linewidth}}
\toprule
原子序数 & 符号 & 电子构型 & 单电子数 \\
\midrule
\endfirsthead
\toprule
原子序数 & 符号 & 电子构型 & 单电子数 \\
\midrule
\endhead
$1$ & $H$ & $1s^1$ & $1$ \\
$2$ & $He$ & $1s^2$ & $0$ \\
$3$ & $Li$ & $[He]2s^1$ & $1$ \\
$4$ & $Be$ & $[He]2s^2$ & $0$ \\
$5$ & $B$ & $[He]2s^22p^1$ & $1$ \\
$6$ & $C$ & $[He]2s^22p^2$ & $2$ \\
$7$ & $N$ & $[He]2s^22p^3$ & $3$ \\
$8$ & $O$ & $[He]2s^22p^4$ & $2$ \\
$9$ & $F$ & $[He]2s^22p^5$ & $1$ \\
$10$ & $Ne$ & $[He]2s^22p^6$ & $0$ \\
$11$ & $Na$ & $[Ne]3s^1$ & $1$ \\
$12$ & $Mg$ & $[Ne]3s^2$ & $0$ \\
$13$ & $Al$ & $[Ne]3s^23p^1$ & $1$ \\
$14$ & $Si$ & $[Ne]3s^23p^2$ & $2$ \\
$15$ & $P$ & $[Ne]3s^23p^3$ & $3$ \\
$16$ & $S$ & $[Ne]3s^23p^4$ & $2$ \\
$17$ & $Cl$ & $[Ne]3s^23p^5$ & $1$ \\
$18$ & $Ar$ & $[Ne]3s^23p^6$ & $0$ \\
$19$ & $K$ & $[Ar]4s^1$ & $1$ \\
$20$ & $Ca$ & $[Ar]4s^2$ & $0$ \\
$21$ & $Sc$ & $[Ar]3d^14s^2$ & $1$ \\
$22$ & $Ti$ & $[Ar]3d^24s^2$ & $2$ \\
$23$ & $V$ & $[Ar]3d^34s^2$ & $3$ \\
$24$ & $Cr$ & $[Ar]3d^54s^1$ & $6$ \\
$25$ & $Mn$ & $[Ar]3d^54s^2$ & $5$ \\
$26$ & $Fe$ & $[Ar]3d^64s^2$ & $4$ \\
$27$ & $Co$ & $[Ar]3d^74s^2$ & $3$ \\
$28$ & $Ni$ & $[Ar]3d^84s^2$ & $2$ \\
$29$ & $Cu$ & $[Ar]3d^{10}4s^1$ & $1$ \\
$30$ & $Zn$ & $[Ar]3d^{10}4s^2$ & $0$ \\
$31$ & $Ga$ & $[Ar]3d^{10}4s^24p^1$ & $1$ \\
$32$ & $Ge$ & $[Ar]3d^{10}4s^24p^2$ & $2$ \\
$33$ & $As$ & $[Ar]3d^{10}4s^24p^3$ & $3$ \\
$34$ & $Se$ & $[Ar]3d^{10}4s^24p^4$ & $2$ \\
$35$ & $Br$ & $[Ar]3d^{10}4s^24p^5$ & $1$ \\
$36$ & $Kr$ & $[Ar]3d^{10}4s^24p^6$ & $0$ \\
\bottomrule
\end{longtable}
\endgroup

\section{元素性质递变规律}
\subsection{元素周期表中区的划分}
\begin{itemize}
\item s 区：IA、IIA，$ns^{1\sim2}$，$H$ 和较活泼金属，属于主族
\end{itemize}

\begin{itemize}
\item d 区：IIIB $\sim$ VIIB 和 VIII，$(n-1)d^{1\sim9}ns^{1\sim2}$，全部是金属元素
\end{itemize}

\begin{itemize}
\item ds 区：IB、IIB，$(n-1)d^{10}ns^{1\sim2}$,全是金属元素
\end{itemize}

\begin{itemize}
\item p 区：IIIA $\sim$ VIIA 和第 $0$ 族，$ns^2np^{1\sim6}$，除 $H$ 外的所有非金属，属于主族
\end{itemize}

\begin{itemize}
\item f 区：镧系、锕系
\end{itemize}
\subsection{价电子构型}
主族元素：价电子为最外层电子\par

副族元素：价电子为最外层和次外层电子\par
\section{元素周期律（下）}
以下规律仅限主族元素\par

\subsection{第一电离能排布呈周期性变化}
定义：基态气态原子失去最外层 $X$ 个电子变为 $+X$ 价离子所需的最小能量成为第 $X$ 电离能\par

规律：\par

\begin{itemize}
\item 同周期从左到右，第一电离能从小到大排序为 $1\mbox{、}3\mbox{、}2\mbox{、}4\mbox{、}6\mbox{、}5\mbox{、}7\mbox{、}8$
\end{itemize}

\begin{itemize}
\item 同族从上到下，第一电离能依次减小
\end{itemize}

\begin{itemize}
\item 若一个元素第 $X$ 电离能和第 $(X+1)$ 电离能相对差距较大，则大概率为 XA 族元素
\end{itemize}

\subsection{亲和能呈周期性变化}
定义：原子得到 $X$ 个电子变为 $-X$ 价离子所需的最小能量成为第 $X$ 亲和能\par

$^*$规律：同周期从左到右大致减小\par

$^*$因反例太多，故只是大致规律\par
\subsection{电负性呈周期性变化}
定义：“定量”的非金属性，与原子吸引电子能力正相关\par

规律（稀有元素气体不考虑）：\par

\begin{itemize}
\item 同周期从左到右，电负性依次增加
\end{itemize}

\begin{itemize}
\item 同族从上到下，电负性依次减小
\end{itemize}

\begin{itemize}
\item 电负性越大，非金属性越强
\end{itemize}

意义：\par

1. 判断元素是否为金属\par

通常电负性 $<1.8$ 是金属，$>1.8$ 是非金属，$>2.6$ 是典型非金属\par

2. 判断化合价正负\par

共价键化合物中，电负性大者为负价，小者为正价（eg. $H$ 电负性为 $2.1$，比金属、$Si$、$B$ 大，故和它们的化合物中呈 $-1$ 价，其余呈 $+1$ 价）\par

3. 判断化合物类型（仅限二元化合物）\par

电负性之差 $>1.7$，多为离子化合物（例外：$HF$）；\par

电负性之差 $<1.7$，多为共价化合物（例外：$KH$、$NaH$）\par

共价化合物中差值越大，极性越强（相同元素间电负性之差为 $0$，不影响化合价）\par
\section{金属键与金属晶体}
\subsection{金属键}
金属离子与自由电子之间强烈的相互作用称为金属键\par

原子化热：$1$ mol 金属固体完全气化成相互远离的气态原子时吸收的能量\par
\subsection{金属键的性质}
金属元素原子半径越小、单位体积内自由电子数目越多，金属键越强（硬度越大、熔沸点越高）\par

金属键与金属性强弱大致负相关\par
\subsection{金属晶体堆积方式}
简单立方（eg. $Po$）、体心立方（eg. $Na$、$K$、$Cr$、$Mo$、$W$）、面心立方（eg. $Au$、$Ag$、$Cu$、$Pb$）\par

简单立方：正方体 $8$ 个顶点各有一个原子，一个晶胞的原子数：$8\times \dfrac{1}{8}=1$\par

体心立方：正方体 $8$ 个顶点和正方体中心各有一个原子，一个晶胞的原子数：$8\times \dfrac{1}{8}+1=2$\par

面心立方：正方体 $8$ 个顶点和 $6$ 个面中心各有一个原子，一个晶胞的原子数：$8\times \dfrac{1}{8}+6\times \dfrac{1}{2}=4$\par
\subsection{相关计算}
\subsubsection{简单立方：棱相切}
$a=2r$，$\rho=\dfrac{m}{V}=\dfrac{\frac{M}{N_A}}{(2r)^3}$，空间利用率 $\dfrac{\frac{4}{3}\pi r^3}{a^3}\approx52.3\%$，$6$ 配位\par

\subsubsection{体心立方：体对角线相切}
$a=\dfrac{4\sqrt{3}}{3}r$，空间利用率 $2\times \dfrac{\frac{4}{3}\pi r^3}{a^3}\approx68.0\%$，$8$ 配位\par

\subsubsection{面心立方：面对角线相切}
$a=2\sqrt{2}r$，空间利用率 $4\times \dfrac{\frac{4}{3}\pi r^3}{a^3}\approx74\%$（是同一元素原子的最大利用率），$12$ 配位\par
\section{离子键与离子晶体}
\subsection{离子键}
阴阳离子的静电作用成为离子键\par

离子键作用力大（$MgO>NaF>NaCl>KCl>KBr$），熔沸点高，比热容、硬度大，一般溶解度大\par
\subsection{离子晶体}
由阴阳离子按一定方式有规则地排列形成的晶体成为离子晶体（包括多数活泼金属化合物、$NH_4^+$ 化合物和 $Al_2O_3$）\par

$1$ mol 固体离子化合物变为气态正负离子吸收的能量称为晶格能\par
\subsection{氯化钠形结构}
$Na^+$ 周围 $6$ 个 $Cl^-$，$Cl^-$ 周围 $6$ 个 $Na^+$（上下前后左右）\par

$Na^+$ 周围 $12$ 个 $Na^+$，$Cl^-$ 周围 $12$ 个 $Cl^-$（对角线）\par

每个 $NaCl$ 晶胞中有 $4$ 个 $Na^+$ 和 $4$ 个 $Cl^-$\par

常见的氯化钠形离子晶体还有 $KCl$、$NaBr$、$CaO$、$MgO$\par
\subsection{氯化铯形结构}
$Cs^+$ 周围 $8$ 个 $Cl^-$，$Cl^-$ 周围 $8$ 个 $Cs^+$\par

$Cl^-$ 周围 $6$ 个 $Cl^-$\par

每个 $CsCl$ 晶胞中有 $1$ 个 $Cl^-$ 和 $1$ 个 $Cs^+$\par

常见的氯化铯形晶体还有 $CsBr$、$CsI$、$NH_4Cl$\par
\section{共价键与共价晶体}
\subsection{共价键}
通常情况下，吸引电子能力相近（即电负性相差不大）的原子之间通过共用电子对形成共价键\par

（1）共价键的本质：成键原子相互靠近时，原子轨道发生重叠，自旋方向相反的未成对电子形成共用电子对，两原子核间的电子云密度增加，体系能量降低\par

（2）共价键的性质：共价键具有方向性、饱和性\par

（3）共价键的分类\par

\begin{itemize}
\item 单键（$H-H$）、双键（$O=C=O$）、三键（$N\equiv N$）
\end{itemize}

\begin{itemize}
\item 极性键（成键原子不同）、非极性键（成键原子相同）
\end{itemize}

eg. $CH_3-CH_2-OH$ 中，$C-C$ 为非极性键，其余为极性键\par

\begin{itemize}
\item 常见共价键、配位键（eg. $H_3O^+$、$NH_4^+$）
\end{itemize}

\begin{itemize}
\item $\sigma$ 键（头碰头）、$\pi$ 键（肩并肩），大 $\pi$ 键为多原子共同形成的 $\pi$ 键
\end{itemize}

$s$ 轨道只能形成 $\sigma$ 键；双键一般是一个 $\sigma$ 键和一个 $\pi$ 键；三键一般是一个 $\sigma$ 键和两个 $\pi$ 键\par

$\sigma$ 键一般比 $\pi$ 键牢固，故化学反应中 $\pi$ 键较容易断裂\par

（4）共价键的参数\par

\begin{itemize}
\item 键长：两成键原子核间的平均距离
\end{itemize}

\begin{itemize}
\item 键能：在 $101$kPa、$298$K 条件下，$1$mol 气态 $AB$ 分子生成 $1$mol 气态 $A$ 原子和 $1$mol 气态 $B$ 原子的过程中吸收的能量，称为 $AB$ 间共价键的键能
\end{itemize}

\begin{notebox}
注：注意区分原子化热、晶格能、键能\par
\end{notebox}
\begin{itemize}
\item 键角：两共价键间的夹角
\end{itemize}

（5）反应热（$\Delta H$）\par

$\Delta H=$ 生成物能量 $-$ 反应物能量\par

$\Delta H=$ 反应物总键能 $-$ 生成物总键能\par

$\Delta H=$ 断键吸收能量 $-$ 成键放出能量\par

$\Delta H>0$ 为吸热反应，$\Delta H<0$ 为放热反应；$\Delta H$ 越小，放热越多\par
\subsection{共价晶体}
所有原子通过共价键结合，形成空间网状结构的晶体叫做共价晶体\par

共价键键能越大、键长越小，熔沸点越高、硬度越大\par

共价晶体中只有共价键，且一般不溶于水\par

$Si$、$SiC$、$SiO_2$、$AlN$、金刚石都是共价晶体\par
\section{分子间作用力与分子晶体}
\subsection{范德华力}
范德华力没有饱和性、方向性；对于组成、结构相似的分子，其范德华力一般随相对分子质量的增大而增大\par
\subsection{氢键}
氢键一般是 $H$ 与 $N/O/F$ 相连时形成的分子间或分子内作用力\par

氢键形成的条件：有孤电子对、有 $H$、有 $N/O/F$\par
\subsection{分子晶体}
分子晶体是分子通过分子间作用力构成的固态物质，它们一般熔点低、硬度小\par

以共价键结合的物质中，多数非金属单质（$H_2$、$Cl_2$、$N_2$）和非金属元素组成的无机化合物（$CO_2$、$HCl$、$H_2O$）都属于分子晶体\par
\section{分子的空间结构}
\subsection{杂化轨道理论}
（1）$sp^3$ 杂化，键角约 $109^{\circ}28'$，正四面体（eg. $CH_4$）\par

（2）$sp^2$ 杂化，键角约 $120^{\circ}$，平面正三角形（eg. $BCl_3$）\par

（3）$sp$ 杂化，键角 $180^{\circ}$，直线形（eg. $BeCl_2$）\par

（4）$sp^3d$ 杂化，三角双锥形（eg. $PCl_5$）\par

（5）$sp^3d^2$ 杂化，正八面体（eg. $SF_6$）\par

孤电子对数 $=\dfrac{1}{2}($价层电子数 $-$ 成键电子数$)$\par

价电子对数 $=$ 孤电子对数 $+$ $\sigma$ 键数 $=$ 参与杂化的轨道数\par
\subsection{价层电子对互斥理论}
（1）价电子对数的计算\par

eg. $NO_3^-$：$\dfrac{1}{2}(5+1-3\times2)+3=3$\par

$PCl_3$：$\dfrac{1}{2}(5-3\times1)+3=4$\par

$SO_3$：$\dfrac{1}{2}(6-3\times2)+3=3$\par

$SO_2$：$\dfrac{1}{2}(6-2\times2)+2=3$\par

（2）分类\par

\begin{itemize}
\item $4$ 对：$sp^3$ 杂化，VSEPR 模型为正四面体
\end{itemize}

$4$ 键：正四面体（eg. $CH_4$、$CCl_4$、$SiH_4$）\par

$3$ 键：三角锥形（eg. $NH_3$、$PH_3$、$PCl_3$）\par

$2$ 键：$V$ 形（eg. $H_2O$、$H_2S$）\par

\begin{itemize}
\item $3$ 对：$sp^2$ 杂化，VSEPR 模型为平面正三角形
\end{itemize}

$3$ 键：平面正三角形（eg. $SO_3$、$NO_3^-$、$CO_3^{2-}$）\par

$2$ 键：$V$ 形（eg. $SO_2$、$NO_2^-$）\par

\begin{itemize}
\item $2$ 对：$sp$ 杂化，直线形（eg. $CO_2$、$BeCl_2$）
\end{itemize}

\begin{notebox}
扩展：$XeF_4$ 的价电子对数为 $\dfrac{1}{2}(8-4\times1)+4=6$，VSEPR 模型为正八面体，分子构型为平面正方形\par
\end{notebox}
\subsection{等电子原理}
具有相同原子数和价电子数的 $AB_n$ 型分子或离子具有相同的结构特征（但键角不一定相同）\par

eg. $CH_4$：$\dfrac{1}{2}(4-4\times1)+4=4$；$NH_4^+$：$\dfrac{1}{2}(5-1-4\times1)+4=4$，故 $CH_4$ 和 $NH_4^+$ 构型相同\par
\section{分子的极性}
\subsection{定义}
正负电荷重心不重合的分子为极性分子，正负电荷重心重合的分子为非极性分子\par
\subsection{非极性分子分类}
（1）单质分子（$O_3$ 除外）\par

（2）完全对称的 $AB_n$ 型分子\par

eg. $CH_4$、$CCl_4$、$PCl_5$、$SF_6$、$IF_7$、$BF_3$、$BCl_3$、$SO_3$、$CO_2$\par
\subsection{相似相溶}
极性溶质易溶于极性溶剂，非极性溶质易溶于非极性溶剂\par

离子化合物一般易溶于水（除晶格能特别大的物质外）\par
\section{配合物}
\subsection{各部分名称}
$[Zn(NH_3)_4]SO_4$ 中，$[Zn(NH_3)_4]^{2+}$ 是内界，$Zn^{2+}$ 是中心原子，$NH_3$ 是配体，$4$ 是配位数，$SO_4^{2-}$ 是外界，名称为“硫酸四氨合锌(II)”\par

若内界不显电性，则无外界，命名时不加中括号\par

eg. $[Cu(NH_3)_2(H_2O)_2]SO_4$：硫酸二氨二水合铜(II)\par

$K_3[Fe(SCN)_6]$：六硫氰合铁(III)酸钾\par
\subsection{配位键}
由一个原子（或离子）单独提供电子对形成的共价键称为配位键\par
\subsection{形成方式}
eg. $[Cu(NH_3)_4]^{2+}$ 为深蓝色\par

$Cu^{2+}$ 的价电子构型可看作 $3d^9$，采取 $dsp^2$ 杂化，形成四个空轨道，因此 $[Cu(NH_3)_4]^{2+}$ 常按平面正方形结构处理\par
\subsection{相关性质}
（1）配位数一般为中心原子化合价的 $2$ 倍，也有可能是 $4$ 或 $6$ 倍\par

（2）中心原子一般为金属\par

（3）配体一般为非金属化合物或离子\par

（4）配位能力与电负性大致负相关，电负性越大，配位能力越弱\par
\chapter{化学反应原理}
\section{化学反应的热效应}
\subsection{反应热}
（1）定义：在等温条件下，化学反应体系向环境释放或从环境吸收的能量\par

（2）定性计算\par

$\Delta H=$ 生成物能量 $-$ 反应物能量\par

$\Delta H=$ 反应物总键能 $-$ 生成物总键能\par

$\Delta H=$ 正反应活化能 $-$ 逆反应活化能\par

$\Delta H<0$ 为放热反应，$\Delta H>0$ 为吸热反应\par
\subsection{热化学方程式}
（1）定义：表明反应所释放或吸收热量的方程式\par

（2）与一般方程式的不同点\par

\begin{itemize}
\item 需要注明各物质状态
\end{itemize}

\begin{itemize}
\item 需要注明反应热 $\Delta H$
\end{itemize}

\begin{itemize}
\item 化学计量数可以为分数，且与 $\Delta H$ 正相关
\end{itemize}
\subsection{反应热的测量}
仪器：（简易）量热计\par

步骤：\par

1. 量取 $50$mL $0.50mol\cdot L^{-1}$ $HCl$，倒入简易量热计，记录 $HCl$ 温度 $T_1$\par

2. 量取 $50$mL $0.55mol\cdot L^{-1}$ $NaOH$，记录 $NaOH$ 温度 $T_2$\par

3. 将 $NaOH$ 迅速倒入量热计，立即盖上盖板，用环形玻璃搅拌棒轻轻搅拌，记录体系最高温度 $T_3$\par

4. 计算\par

$\Delta T=T_3-\dfrac{1}{2}(T_1+T_2)$\par

$C=(V_{HCl}\rho_{HCl}+V_{NaOH}\rho_{NaOH})\times4.18\ J\cdot ^{\circ}C^{-1}$\par

$\Delta H=-\dfrac{C\cdot\Delta T\cdot10^{-3}}{0.025}\ kJ\cdot mol^{-1}$\par
\subsection{盖斯定律}
一个化学反应，不论是一步完成，还是分几步完成，其总的热效应是完全相等的\par
\subsection{燃烧热与热值}
标准燃烧热：$1$mol 物质完全燃烧的反应热\par

热值：$1$g 物质完全燃烧的反应热\par
\section{电化学}
\subsection{原电池}
负极：$Zn-2e^-=\!=\!=Zn^{2+}$\par

正极：$Cu^{2+}+2e^-=\!=\!=Cu$ 或 $2H^++2e^-=\!=\!=H_2\uparrow$\par

电池反应：$Zn+Cu^{2+}=\!=\!=Zn^{2+}+Cu$ 或 $Zn+2H^+=\!=\!=Zn^{2+}+H_2\uparrow$\par

缺点：若直接将 $Zn$ 放入酸中，$Zn+H_2SO_4=\!=\!=ZnSO_4+H_2\uparrow$，会有内耗\par

改进：使用盐桥，$Cl^-$ 等阴离子向负极移动，$K^+$ 等阳离子向正极移动\par

自发的氧化还原反应一般都能设计成原电池（歧化反应除外）\par
\subsection{电极方程式}
原电池中，失电子、作还原剂的是负极；得电子、作氧化剂的是正极\par

电解池中，发生氧化反应的是阳极，发生还原反应的是阴极\par

eg. 铅蓄电池放电\par

负极：$Pb+SO_4^{2-}-2e^-=\!=\!=PbSO_4$\par

正极：$PbO_2+SO_4^{2-}+4H^++2e^-=\!=\!=PbSO_4+2H_2O$\par

电池反应：$Pb+PbO_2+2SO_4^{2-}+4H^+=\!=\!=2PbSO_4+2H_2O$\par

\medskip
eg. 甲烷-氧气燃料电池（碱性）\par

负极：$CH_4+10OH^--8e^-=\!=\!=CO_3^{2-}+7H_2O$\par

正极：$O_2+2H_2O+4e^-=\!=\!=4OH^-$\par

\medskip
eg. 银锌纽扣电池\par

负极：$Zn+2OH^--2e^-=\!=\!=ZnO+H_2O$\par

正极：$Ag_2O+H_2O+2e^-=\!=\!=2Ag+2OH^-$\par

\medskip
eg. 氢氧燃料电池\par

碱性：\par

负极：$H_2-2e^-+2OH^-=\!=\!=2H_2O$\par

正极：$O_2+4e^-+2H_2O=\!=\!=4OH^-$\par

酸性：\par

负极：$H_2-2e^-=\!=\!=2H^+$\par

正极：$O_2+4H^++4e^-=\!=\!=2H_2O$\par

\medskip
eg. $CO-O_2$ 燃料电池（熔融 $Na_2CO_3$）\par

负极：$CO-2e^-+CO_3^{2-}=\!=\!=2CO_2$\par

正极：$O_2+2CO_2+4e^-=\!=\!=2CO_3^{2-}$\par

\medskip
eg. 锌锰电池\par

负极：$Zn+2OH^--2e^-=\!=\!=Zn(OH)_2$\par

正极：$MnO_2+H_2O+e^-=\!=\!=MnOOH+OH^-$\par
\subsection{电解池}
电解池中，与电源正极相连的是阳极，发生氧化反应；与电源负极相连的是阴极，发生还原反应\par

惰性电极（eg. $Au$、$Pt$、石墨）时，常见放电顺序：\par

阳极：$S^{2-}>I^->Br^->Cl^->OH^->F^-$、$SO_4^{2-}$、$NO_3^-$、$CO_3^{2-}$\par

阴极：$Ag^+>Fe^{3+}>Cu^{2+}>H^+$（大量）$>Fe^{2+}>Zn^{2+}>H^+$（水中）$>Al^{3+}>Mg^{2+}>Na^+>Ca^{2+}$\par

eg. 电解 $NaCl(aq)$\par

阳极：$2Cl^--2e^-=\!=\!=Cl_2\uparrow$\par

阴极：$2H_2O+2e^-=\!=\!=H_2\uparrow+2OH^-$\par

\medskip
eg. 电解 $K_2MnO_4(aq)$\par

阳极：$MnO_4^{2-}-e^-=\!=\!=MnO_4^-$\par

阴极：$2H_2O+2e^-=\!=\!=H_2\uparrow+2OH^-$\par

\medskip
eg. 电解 $CuSO_4(aq)$\par

阳极：$2H_2O-4e^-=\!=\!=4H^++O_2\uparrow$\par

阴极：$Cu^{2+}+2e^-=\!=\!=Cu\downarrow$\par

\medskip
eg. 电解 $NaHCO_3(aq)$ 生成 $CH_4$\par

$10HCO_3^-+8e^-=\!=\!=CH_4\uparrow+3H_2O+9CO_3^{2-}$\par
\subsection{电镀}
eg. 为导电塑料镀铜\par

阴极：导电塑料（镀件），阳极：$Cu$（镀层金属），电解液：$CuSO_4(aq)$\par

阳极：$Cu-2e^-=\!=\!=Cu^{2+}$\par

阴极：$Cu^{2+}+2e^-=\!=\!=Cu\downarrow$\par
\subsection{精炼}
eg. 铜的电解精炼\par

阳极：粗铜，阴极：纯铜，电解液：$CuSO_4(aq)$\par

阳极：$Zn-2e^-=\!=\!=Zn^{2+}$，\\[0.25em]
\hspace*{2em}$Fe-2e^-=\!=\!=Fe^{2+}$，\\[0.25em]
\hspace*{2em}$Cu-2e^-=\!=\!=Cu^{2+}$\par

阴极：$Cu^{2+}+2e^-=\!=\!=Cu\downarrow$\par

阳极泥：主要成分为 $Au$、$Ag$、$Pt$\par
\subsection{电化学腐蚀}
（1）析氢腐蚀（生成 $H_2$）\par

eg. 铁表面为酸性水膜的析氢腐蚀\par

负极：$Fe-2e^-=\!=\!=Fe^{2+}$\par

正极：$2H^++2e^-=\!=\!=H_2\uparrow$\par

总反应：$Fe+2H^+=\!=\!=Fe^{2+}+H_2\uparrow$\par

（2）吸氧腐蚀（吸收 $O_2$）\par

eg. 铁在潮湿空气中的吸氧腐蚀\par

负极：$Fe-2e^-=\!=\!=Fe^{2+}$\par

正极：$2H_2O+O_2+4e^-=\!=\!=4OH^-$\par

总反应：$2Fe+2H_2O+O_2=\!=\!=2Fe(OH)_2$\par

$4Fe(OH)_2+2H_2O+O_2=\!=\!=4Fe(OH)_3$，最终形成 $Fe_2O_3\cdot nH_2O$\par

铁锈是红褐色的\par
\subsection{电化学防护}
（1）刷油漆、涂一层膜\par

（2）改变结构（eg. 不锈钢）\par

（3）牺牲阳极法：把被保护金属与更活泼金属连接\par

（4）外接电流法（阴极保护法）\par
\section{化学反应速率}
\subsection{碰撞理论}
微粒要有有效碰撞才能反应\par

因此，发生反应需要：\par

1. 碰撞\par

2. 在一定方向上碰撞\par

3. 有一定能量\par

活化分子数、碰撞次数影响单位时间内有效碰撞次数，即化学反应速率\par
\subsection{化学反应速率的影响因素}
以下规律主要适用于基元反应\par

（1）温度：活化分子数增大、碰撞次数增多，一般升高 $10^{\circ}C$，速率变为原来的 $2\sim4$ 倍\par

\begin{notebox}
扩展：$2NO+O_2=\!=\!=2NO_2$ 中，温度越高，反应反而越慢\par
\end{notebox}
（2）浓度：浓度增大，碰撞次数增多；不适用于固体和纯液体\par

（3）催化剂：改变反应历程，降低活化能，活化分子数增多\par

（4）压强：对气体反应，增大压强相当于增大浓度；影响反应的是分压，不是总压强\par

（5）接触面积\par

（6）原电池\par
\subsection{定量表示}
$v=\dfrac{\Delta c}{\Delta t}$，常用单位为 $mol\cdot L^{-1}\cdot s^{-1}$\par

eg. $N_2+3H_2\rightleftharpoons2NH_3$（$10$L 容器）\par

$0$s 时：$n(N_2)=10$mol，$n(H_2)=30$mol，$n(NH_3)=0$mol\par

$10$s 时：$n(N_2)=2$mol，$n(H_2)=6$mol，$n(NH_3)=16$mol\par

$\Delta n(N_2)=8$mol，$\Delta n(H_2)=24$mol，$\Delta n(NH_3)=16$mol\par

$v(N_2)=0.08mol\cdot L^{-1}\cdot s^{-1}$\par

$v(H_2)=0.24mol\cdot L^{-1}\cdot s^{-1}$\par

$v(NH_3)=0.16mol\cdot L^{-1}\cdot s^{-1}$\par

反应速率要说明是哪种物质；同一反应中，同一时间段内各物质反应速率之比等于化学计量数之比\par
\subsection{化学反应速率的决定因素}
化学反应速率的决定因素是反应物本身\par
\section{化学平衡}
\subsection{前提}
可逆反应在一定条件下才能建立化学平衡\par
\subsection{表现}
各物质的量、浓度、百分含量不变\par
\subsection{本质}
$v_{\scriptscriptstyle\mbox{正}}=v_{\scriptscriptstyle\mbox{逆}}$\par

如果条件改变，平衡可能被打破\par
\subsection{表示}
（1）压强不变（反应物、生成物气体的量不同时）\par

（2）质量不变（反应物、生成物中有固体时）\par

（3）颜色不变（反应有颜色变化时）\par
\subsection{平衡的打破}
（1）浓度\par

增大反应物浓度或减小生成物浓度，平衡右移；减小反应物浓度或增大生成物浓度，平衡左移\par

（2）压强（反应物、生成物中有气体时）\par

根据浓度商 $Q$ 与平衡常数 $K$ 判断方向。eg. $N_2+3H_2\rightleftharpoons2NH_3$，$K=\dfrac{c^2(NH_3)}{c(N_2)\cdot c^3(H_2)}$，$K$ 只与温度有关\par

$Q<K$ 时右移，$v_{\scriptscriptstyle\mbox{正}}>v_{\scriptscriptstyle\mbox{逆}}$；$Q=K$ 时平衡，$v_{\scriptscriptstyle\mbox{正}}=v_{\scriptscriptstyle\mbox{逆}}$；$Q>K$ 时左移，$v_{\scriptscriptstyle\mbox{正}}<v_{\scriptscriptstyle\mbox{逆}}$\par

（3）温度\par

升温向 $\Delta H>0$ 的吸热方向移动，降温向 $\Delta H<0$ 的放热方向移动\par

eg. $2NO_2(g)\rightleftharpoons N_2O_4(g)$，$\Delta H<0$，降温平衡右移\par

（4）催化剂\par

催化剂不影响平衡，只同时改变正、逆反应速率\par
\subsection{勒夏特列原理}
改变影响化学平衡的一个因素，平衡将向着能够减弱这种改变的方向移动\par
\section{电离平衡}
\subsection{强弱电解质}
水溶液中完全电离的电解质是强电解质（强酸、强碱、大多数盐）\par

水溶液中部分电离的电解质是弱电解质（弱酸、弱碱、$HgCl_2$、$Pb(Ac)_2$ 等）\par
\subsection{判断 $HA$ 是强酸还是弱酸}
（1）$0.1mol\cdot L^{-1}$ $HA$，$pH=1$ 为强酸，$pH>1$ 为弱酸\par

（2）和同浓度 $HCl$ 比导电性，相同为强，更弱为弱\par

（3）和同浓度 $HCl$ 与 $Zn$ 反应，产生 $H_2$ 速率相同为强，更慢为弱\par

（4）升高温度，$pH$ 变小为弱酸，$pH$ 基本不变为强酸\par

（5）稀释 $100$ 倍，$\Delta pH=2$ 为强酸，$\Delta pH<2$ 为弱酸\par

（6）加 $NaA(s)$，$pH$ 升高为弱酸，$pH$ 基本不变为强酸\par

（7）计算电离常数 $K_a=\dfrac{c(H^+)\cdot c(A^-)}{c(HA)}$，$K_a$ 越大，酸性越强\par
\subsection{水的电离}
$H_2O\rightleftharpoons H^++OH^-$，吸热\par

\begin{notebox}
注：几乎所有常见弱电解质电离都吸热\par
\end{notebox}
$K_w=c(H^+)\cdot c(OH^-)=1.0\times10^{-14}$（$25^{\circ}C$），温度升高，$K_w$ 增大\par
\subsection{溶液酸碱性}
$c(H^+)>c(OH^-)$ 为酸性；$c(H^+)=c(OH^-)$ 为中性；$c(H^+)<c(OH^-)$ 为碱性\par

$25^{\circ}C$ 时，$c(H^+)>10^{-7}$ 为酸性，$c(H^+)=10^{-7}$ 为中性，$c(H^+)<10^{-7}$ 为碱性\par
\subsection{$pH$ 值}
$pH=-\lg c(H^+)$\par

eg. $pH=6$ 的 $HCl(aq)$ 稀释 $100$ 倍，求 $pH$\par

稀释后 $c_{HCl}=10^{-8}mol\cdot L^{-1}$，设水电离产生的 $c(OH^-)=x$，则\par

$(x+10^{-8})x=10^{-14}$，$x\approx9.5\times10^{-8}mol\cdot L^{-1}$\par

$c(H^+)=10^{-8}+9.5\times10^{-8}=1.05\times10^{-7}mol\cdot L^{-1}$\par

$pH=-\lg c(H^+)\approx6.98$\par
\subsection{中和滴定}
（1）原理：$H^++OH^-=\!=\!=H_2O$，$c_1V_1=c_2V_2$\par

（2）指示剂\par

酚酞（$8.1\sim10$）：无色 $\to$ 浅红 $\to$ 红\par

甲基橙（$3.1\sim4.4$）：红 $\to$ 橙 $\to$ 黄\par

石蕊（不常用）：红 $\to$ 紫 $\to$ 蓝\par

酸滴定碱：甲基橙（黄 $\to$ 橙）\par

碱滴定酸：酚酞（无色 $\to$ 浅红）\par

（3）滴定管\par

酸式滴定管可盛放酸性溶液、有机溶剂、氧化剂；碱式滴定管可盛放碱性溶液\par

（4）步骤\par

1. 洗滴定管：自来水 $\to$ 蒸馏水 $\to$ 待装液润洗三次\par

2. 洗锥形瓶：用蒸馏水洗三次\par

3. 滴定管装液体，装至 $0$ 刻度或以下\par

4. 滴定：已知浓度测未知浓度，平行滴定三次，取平均值\par

5. 记录并处理数据\par

（5）终点现象\par

碱滴定酸：滴入最后半滴碱溶液，溶液由无色变为浅红色，且半分钟内不褪色\par

酸滴定碱：滴入最后半滴酸溶液，溶液由黄色变为橙色，且半分钟内不褪色\par
\section{水解平衡}
\subsection{水解产生原因}
eg. $NaAc(aq)$\par

$HAc\rightleftharpoons H^++Ac^-$，\\[0.25em]
\hspace*{2em}$H_2O\rightleftharpoons H^++OH^-$，\\[0.25em]
\hspace*{2em}$NaAc=\!=\!=Na^++Ac^-$\par

$Ac^-+H_2O\rightleftharpoons HAc+OH^-$，故 $NaAc(aq)$ 显碱性\par
\subsection{水解平衡常数 $K_h$}
$HAc\rightleftharpoons H^++Ac^-$，$K_a=\dfrac{c(H^+)\cdot c(Ac^-)}{c(HAc)}$\par

$H_2O\rightleftharpoons H^++OH^-$，$K_w=c(H^+)\cdot c(OH^-)$\par

$Ac^-+H_2O\rightleftharpoons HAc+OH^-$，$K_h=\dfrac{c(HAc)\cdot c(OH^-)}{c(Ac^-)}=\dfrac{K_w}{K_a}$\par
\subsection{水解规律}
有弱才水解，越弱越水解，谁强显谁性\par

强碱弱酸盐显碱性；强酸弱碱盐显酸性；强酸强碱盐显中性；弱酸弱碱盐按 $K_a$、$K_b$ 比较\par
\subsection{影响水解的因素}
（1）$K_a/K_b$\par

（2）浓度：$c$ 增大，水解程度 $\alpha$ 减小\par

（3）温度：水解吸热，温度升高，$K_h$ 增大，$\alpha$ 增大\par

（4）酸碱性\par

eg. $Fe^{3+}+3H_2O\rightleftharpoons Fe(OH)_3+3H^+$，加入 $HCl$ 可抑制水解，使平衡左移\par
\subsection{双水解}
eg. $Al^{3+}+3H_2O\rightleftharpoons Al(OH)_3+3H^+$\par

$HCO_3^-+H_2O\rightleftharpoons H_2CO_3+OH^-$\par

$Al^{3+}+3HCO_3^-=\!=\!=Al(OH)_3\downarrow+3CO_2\uparrow$\par

泡沫灭火器反应：$Al_2(SO_4)_3+6NaHCO_3=\!=\!=3Na_2SO_4+2Al(OH)_3+6CO_2\uparrow$\par

弱酸酸根离子和弱碱金属离子大部分可以双水解，但若发生氧化还原，或生成物溶解度较大且仍停留在溶液中，则不一定发生完全双水解\par
\subsection{酸式盐酸碱性}
eg. $HCO_3^-$\par

$HCO_3^-+H_2O\rightleftharpoons H_2CO_3+OH^-$，$K_h\approx10^{-7}$\par

$HCO_3^-\rightleftharpoons H^++CO_3^{2-}$，$K_{a2}\approx10^{-11}$\par

$K_h>K_{a2}$，故水解大于电离，溶液显碱性\par

对二元弱酸的酸式酸根，$K_h=\dfrac{K_w}{K_{a1}}$，可比较 $K_w$ 与 $K_{a1}K_{a2}$：\par

$K_{a1}K_{a2}<10^{-14}$ 时显碱性，$K_{a1}K_{a2}>10^{-14}$ 时显酸性\par
\section{离子浓度}
\subsection{电荷守恒}
溶液中正负电荷数相等\par

eg. $NaAc(aq)$ 中，$c(Na^+)+c(H^+)=c(Ac^-)+c(OH^-)$\par
\subsection{物料守恒}
反应前后同种元素原子个数不变，不同元素原子个数比不变\par

eg. $NaAc(aq)$ 中，$c(Na^+)=c(Ac^-)+c(HAc)$\par

$Na_2CO_3(aq)$ 中，$c(Na^+)=2[c(CO_3^{2-})+c(HCO_3^-)+c(H_2CO_3)]$\par

$HAc$ 与 $NaAc$ $1:1$ 混合液中，$2c(Na^+)=c(Ac^-)+c(HAc)$\par
\subsection{运用}
（1）$HAc(aq)$ 与 $NaAc(aq)$ $1:1$ 混合\par

$2c(Na^+)=c(Ac^-)+c(HAc)$\par

$c(Na^+)+c(H^+)=c(Ac^-)+c(OH^-)$\par

若溶液呈酸性，则 $c(H^+)>c(OH^-)$，所以 $c(Na^+)<c(Ac^-)$\par

故 $c(Ac^-)>c(Na^+)>c(HAc)>c(H^+)>c(OH^-)$\par

\medskip
（2）$NaHCO_3(aq)$\par

$c(Na^+)+c(H^+)=2c(CO_3^{2-})+c(HCO_3^-)+c(OH^-)$\par

$c(Na^+)=c(HCO_3^-)+c(CO_3^{2-})+c(H_2CO_3)$\par

故 $c(Na^+)>c(HCO_3^-)>c(OH^-)>c(H^+)$\par

$c(Na^+)>c(HCO_3^-)>c(H_2CO_3)>c(CO_3^{2-})$\par

\medskip
（3）$Na_2CO_3(aq)$\par

$c(Na^+)=2[c(CO_3^{2-})+c(HCO_3^-)+c(H_2CO_3)]$\par

$c(Na^+)+c(H^+)=2c(CO_3^{2-})+c(HCO_3^-)+c(OH^-)$\par

故 $c(Na^+)>c(CO_3^{2-})>c(HCO_3^-)>c(H_2CO_3)$\par

\medskip
（4）往 $NH_4HSO_4(aq)$ 中加 $Ba(OH)_2(aq)$ 至显中性\par

$c(NH_4^+)+c(H^+)=c(OH^-)+2c(SO_4^{2-})$\par

显中性时 $c(H^+)=c(OH^-)$，所以 $c(NH_4^+)=2c(SO_4^{2-})$\par

\medskip
（5）$HCN(aq)$ 与 $NaCN(aq)$ $1:1$ 混合，$K_a(HCN)\approx10^{-10}$\par

$2c(Na^+)=c(CN^-)+c(HCN)$\par

$c(Na^+)+c(H^+)=c(CN^-)+c(OH^-)$\par

$K_a(HCN)=\dfrac{c(CN^-)\cdot c(H^+)}{c(HCN)}\approx c(H^+)\approx10^{-10}$，故 $c(H^+)<c(OH^-)$，且 $c(Na^+)>c(CN^-)$\par

因此 $c(HCN)>c(Na^+)>c(CN^-)>c(OH^-)>c(H^+)$\par
\section{沉淀溶解平衡（$K_{sp}$）}
\subsection{前提}
沉淀溶解平衡适用于难溶电解质，且通常是固体溶于液体\par
\subsection{实例：$AgCl$ 的溶解}
$AgCl(s)\rightleftharpoons Ag^+(aq)+Cl^-(aq)$\par

\begin{notebox}
注：溶解平衡用 $\rightleftharpoons$，电离用 $=\!=\!=$，因为 $AgCl$ 是强电解质\par
\end{notebox}
$25^{\circ}C$ 时，$K_{sp}=c(Ag^+)\cdot c(Cl^-)=1.8\times10^{-10}$\par

（1）定量\par

$c(Ag^+)=c(Cl^-)=1.34\times10^{-5}mol\cdot L^{-1}$\par

故 $1$L 水中溶解 $AgCl$ 的质量为 $1.34\times10^{-5}\times143.5=1.9\times10^{-3}$g\par

$AgCl$ 的溶解度约为 $1.9\times10^{-4}$g（以 $100$g 水计）\par

（2）定性\par

在 $AgCl$ 饱和溶液中加入 $NaCl(s)$，平衡左移，$AgCl(s)$ 增多，$c(Cl^-)$ 增大，$c(Ag^+)$ 减小，$K_{sp}$ 不变，这称为同离子效应\par

应用：饱和 $NaCl(aq)$ 中加入 $KCl(s)$，有 $NaCl$ 晶体析出\par

（3）温度\par

$AgCl(s)\rightleftharpoons Ag^+(aq)+Cl^-(aq)$ 为吸热过程，温度升高，$K_{sp}$ 增大\par
\subsection{除杂}
$25^{\circ}C$ 时，$Fe(OH)_3$ 的 $K_{sp}=c(Fe^{3+})\cdot c^3(OH^-)=10^{-39}$\par

对 $1mol\cdot L^{-1}$ $FeCl_3$ 溶液，$c(OH^-)=10^{-13}$ 时，$c(H^+)=10^{-1}$，即 $pH=1$ 时 $Fe^{3+}$ 开始沉淀\par

当 $c(Fe^{3+})=10^{-6}mol\cdot L^{-1}$ 时可认为沉淀完全，此时 $c(OH^-)=10^{-11}$，即 $pH=3$\par

$25^{\circ}C$ 时，$Cu(OH)_2$ 的 $K_{sp}=2.2\times10^{-20}=c(Cu^{2+})\cdot c^2(OH^-)$\par

对 $1mol\cdot L^{-1}$ $CuCl_2$ 溶液，$c(OH^-)=10^{-10}$ 时，$pH=4$，即 $pH=4$ 时开始沉淀\par

当 $c(Cu^{2+})=10^{-6}mol\cdot L^{-1}$ 时，$c(OH^-)\approx10^{-7}$，即 $pH=7$ 时沉淀完全\par

所以，溶液中 $Fe^{3+}$、$Cu^{2+}$ 分离时可将 $pH$ 调为 $3\sim4$，此时 $Fe^{3+}$ 沉淀完全，$Cu^{2+}$ 还未开始沉淀\par

另外，加入 $Cu(OH)_2$ 也可分离混合溶液中的 $Fe^{3+}$：\par

$Cu(OH)_2\rightleftharpoons Cu^{2+}+2OH^-$，$c(OH^-)$ 增大，使 $Fe(OH)_3\rightleftharpoons Fe^{3+}+3OH^-$ 左移，$Fe^{3+}$ 沉淀\par
\subsection{沉淀的转化}
（1）$AgCl$ 转化为 $AgI$\par

$K_{sp}(AgCl)=c(Ag^+)\cdot c(Cl^-)=1.8\times10^{-10}$，$c(Ag^+)=1.34\times10^{-5}mol\cdot L^{-1}$\par

$K_{sp}(AgI)=c(Ag^+)\cdot c(I^-)\approx10^{-17}$，且 $c(I^-)\gg10^{-12}$\par

加入 $I^-$ 后，$Ag^+$ 与 $I^-$ 结合形成 $AgI$ 沉淀，$c(Ag^+)$ 减小，$AgCl(s)\rightleftharpoons Ag^+(aq)+Cl^-(aq)$ 右移，故若 $I^-$ 足够多，$AgCl$ 会逐渐转化为 $AgI$\par

（2）分别在 $0.1mol\cdot L^{-1}$ $I^-$、$Cl^-$ 溶液中加入 $Ag^+$\par

$c(Ag^+)=10^{-16}mol\cdot L^{-1}$ 时，$AgI$ 开始沉淀\par

$c(Ag^+)=10^{-9}mol\cdot L^{-1}$ 时，$AgCl$ 开始沉淀\par

同类型、同浓度的两种物质，$K_{sp}$ 小的先沉淀\par

若一个沉淀能完全转化为另一个沉淀，也无法直接判断 $K_{sp}$ 相对大小： 系数比不一定一样、沉淀类型可能不同\par


\backmatter
\cleardoublepage
\thispagestyle{empty}
\begin{titlepage}
\centering
\vspace*{2.0cm}
\rule{0.58\linewidth}{0.8pt}\par
\vspace{1.0cm}
{\zihao{1}\songti 以微观之眼\par}
\vspace{0.35cm}
{\zihao{1}\songti 见万物之变\par}
\vspace{1.0cm}
\rule{0.58\linewidth}{0.8pt}\par
\vfill
{\zihao{5}\rmfamily Reaction \quad Structure \quad Equilibrium\par}
\vspace{0.45cm}
{\zihao{5} 西安交通大学少年班 \quad 江苏省苏州中学预科基地\par}
\vspace*{1.3cm}
\end{titlepage}

\end{document}
